\documentclass[11pt]{article}
\usepackage[margin=1in]{geometry}
\usepackage{amsmath,amssymb,amsthm}
\usepackage{newtxtext,newtxmath}
\usepackage{booktabs,longtable,array}
\usepackage[hidelinks]{hyperref}
\usepackage{enumitem}
\setlist{itemsep=2pt,topsep=4pt}

\newtheorem{proposition}{Proposition}
\newtheorem{lemma}{Lemma}
\newtheorem{corollary}{Corollary}
\theoremstyle{definition}
\newtheorem{definition}{Definition}
\newtheorem{remark}{Remark}

\newcommand{\E}{\mathbb E}
\newcommand{\Ent}{\mathsf E}
\newcommand{\Own}{\mathsf O}
\newcommand{\fr}{\mathfrak r}
\newcommand{\one}{\mathbf 1}
\newcommand{\lab}[1]{\textsc{#1}}

\title{The order slot: beliefs, informed trading and the main comparisons\\[2pt]
\large Proof package P1 (results R2, R3 and R4)}
\author{}
\date{23 October 2026}

\begin{document}
\maketitle

\section*{Scope and status}

This package proves the core of the order-slot manuscript.
\begin{itemize}
\item Lemmas 1--3 (R2): the flow law, the posterior bounds, price sufficiency and the informed investor's
  payoffs under every strategy. Status: \lab{analytical}.
\item Proposition 2 (R3): competition creates competition, for general classes and priors.
  Status: \lab{analytical}.
\item Proposition 3 (R4): the complete equilibrium set along the benchmark path at prior $1/2$.
  Proposition 5 (R4): the strengths with coexistence at every prior. Status of both: \lab{computer-assisted},
  because both use the two Sturm certificates of Lemma~\ref{lem:path}.
\item Proposition 4: the admissible region and the margins table. Status: \lab{analytical}.
\end{itemize}
Decimal values at calibration B depend on $e$. They are registry values with the status
\lab{computer-assisted}, and the results that print them keep their own label.
Every number in this package comes from the scripts in \texttt{scripts/P1\_core/} and appears in
\texttt{registry/P1\_core.csv}. Section~\ref{sec:scripts} lists the scripts and their run logs.

This is the second version of the package. It applies the corrections of the independent check
(memo \texttt{XP1\_check.md}); the memo \texttt{P1\_core.md} lists each change. The notation follows the shared
labels file: $f_j(x)$ is the probability of order $x$ in class $j$.

%==========================================================================================
\section{The trading stage and the equilibrium object}\label{sec:model}

\subsection{Primitives outside the trading stage}

The primitives outside the trading stage are those of the Laplace manuscript, Sections 2.1, 2.2 and 3.1.
\begin{itemize}
\item \emph{Classes and values.} The challenger's class is $j\in\{H,L\}$, with prior $\pi=\Pr(j=H)\in(0,1)$.
  Given its class, its value is $\theta\sim G_j$, with a finite mean. The high class dominates,
  $G_H(v)\le G_L(v)$ for every $v$, and both supports lie above the reserve, $G_H(p^-)=G_L(p^-)=0$.
\item \emph{Incumbent.} The incumbent's value $R$ has a continuous distribution $F$ on $[0,\bar r]$.
  A strengthening of $F$ is a continuous distribution $\tilde F\le F$. The main case is $R\sim U[0,r]$.
\item \emph{Costs.} The challenger's preparation cost is $C\in\{0,c_H\}$, with $\Pr(C=0)=\rho\in(0,1)$ and $c_H>0$.
\item \emph{Sale.} A second-price auction with reserve $p>0$. Prepared bidders bid their values.
\item \emph{Independence.} The pair $(j,\theta)$, the value $R$, the cost $C$ and the draws of the order slot
  are independent.
\end{itemize}

For a known challenger value $\theta$, $t(\theta)$ is expected target proceeds after the challenger prepares,
and $g(\theta)$ is the challenger's expected gross profit. Without entry the proceeds are $t_0=p[1-F(p)]$.
The class objects are
\[
t_j=\E_j[t(\theta)],\quad g_j=\E_j[g(\theta)],\quad \Delta=t_H-t_L,\quad B(\mu)=g_L+\mu(g_H-g_L),
\]
\[
W_H(F)=\E_H\big[F(\theta)\one\{\theta\ge p\}\big].
\]
Here $\Delta$ is the target-payoff spread $\Delta_T$ of the Laplace manuscript, and $W_H(F)$ is the
probability that a prepared class-$H$ challenger wins. When $g_H>g_L$, the preparation threshold is
$\tau=(c_H-g_L)/(g_H-g_L)$, and $B(\mu)\ge c_H$ exactly when $\mu\ge\tau$.

The proofs use three facts from Proposition 1 of the Laplace manuscript (result R1, proved in another package):
\begin{enumerate}[label=(P\arabic*)]
\item $g_H\ge g_L$, $\Delta\ge 0$ and $t_L\ge t_0$ for every $F$;
\item a strengthening $\tilde F\le F$ weakly lowers $g_H$ and $g_L$, so $B_{\tilde F}(\mu)\le B_F(\mu)$ at every $\mu$;
\item a strengthening weakly raises $\Delta$.
\end{enumerate}

\emph{The benchmark.} $G_H$ and $G_L$ are point masses at $h$ and $\ell$, $R\sim U[0,r]$, and $0<p<\ell<r<h$. Then
\[
t_0=p\Big(1-\frac pr\Big),\quad t_H=\frac r2+\frac{p^2}{2r},\quad t_L=\ell-\frac{\ell^2-p^2}{2r},\quad
g_H=h-t_H,\quad g_L=\frac{\ell^2-p^2}{2r},\quad \Delta(r)=\frac{(r-\ell)^2}{2r}.
\]
Also $g_H-g_L=h-r/2-\ell^2/(2r)>0$ for $\ell<r<h$, and
\[
\tau(r)=\frac{2rc_H-\ell^2+p^2}{2rh-r^2-\ell^2}.
\]
The benchmark values are
\[
(h,\ell,p,\rho,c_H,k,\pi)=(10,1,1/2,1/4,6,1/50,1/2),\qquad (r_0,r_1,r_2)=(6/5,3,18/5).
\]
Then $\tau(r)=(48r-3)/[4(20r-r^2-1)]$, with $\tau(r_0)=195/308$, $\tau(r_1)=141/200$ and $\tau(r_2)=4245/5804$.
These formulas agree with equation (4) of the Laplace manuscript at commit \texttt{01bd5d6} and with the
shared objects of the team rules. The script \texttt{p1\_certificates.py} also integrates the sale payoffs
directly over $R$ at eight strengths. Every difference from the closed forms is zero.

\subsection{The order slot}

The trading stage has one public order opportunity.
\begin{enumerate}
\item The order opportunity comes before the challenger's decision.
\item With probability $\lambda\in(0,1)$ the informed investor fills it. The investor sees the class $j$, not
  the value $\theta$, and chooses $q\in\{-1,0,1\}$. A nonzero order costs $k>0$.
\item With probability $1-\lambda$ a liquidity trader fills it. It sells, abstains or buys one unit with
  probabilities $\eta(-1)=\beta$, $\eta(0)=1-2\beta$ and $\eta(1)=\beta$, where $0<\beta\le 1/3$. Its motive is exogenous.
\item No one observes who filled the slot. The identity draw is independent of the class, the costs and $R$.
\item Competitive market makers see the order $X\in\{-1,0,1\}$ and set $P=\E[V_T\mid X]$. The challenger sees
  the price and its own cost, then decides whether to prepare. Prepared bidders learn their values and bid.
\end{enumerate}
The derived objects are
\[
b=(1-\lambda)\beta,\qquad b_0=(1-\lambda)(1-2\beta),\qquad u=\frac{\lambda}{b},
\]
\[
M(\pi)=\frac{\pi(1+u)}{\pi(1+u)+1-\pi},\qquad m(\pi)=\frac{\pi}{\pi+(1-\pi)(1+u)},\qquad
\kappa(\pi)=\min\{m(\pi),1-M(\pi)\}.
\]
The number $b$ is the probability of a liquidity purchase, and also of a liquidity sale. The number $u$ is the
informed presence mass relative to the liquidity mass in one direction, and $1+u=(b+\lambda)/b$ is the largest
likelihood ratio of a nonzero order (Lemma~\ref{lem:flow}(b)). The odds that an order is informed depend on the
strategies: under full orders they are $\pi u$ after a purchase and $(1-\pi)u$ after a sale. At $\pi=1/2$ we write
$M=(1+u)/(2+u)$ and $m=1-M$.

When the investor fills the slot, its order is the order flow, $X=q$. Its profit is $q[V_T-P(q)]-k|q|$.
The arrival probability $\lambda$ multiplies its ex ante payoff and changes no comparison.

\subsection{The equilibrium object}

\begin{definition}[equilibrium]\label{def:eq}
An equilibrium is a tuple $\mathcal E=(\sigma_H,\sigma_L,P,\mu,\zeta,\text{bids})$ with these components:
\begin{enumerate}[label=(\roman*)]
\item class-contingent order distributions $\sigma_j\in\Delta(\{-1,0,1\})$, $j\in\{H,L\}$;
\item a price on every order, $P:\{-1,0,1\}\to\mathbb R$;
\item a posterior on every order, $\mu:\{-1,0,1\}\to[0,1]$, the probability that the class is $H$;
\item a preparation rule: the zero-cost challenger prepares at every price, and the high-cost challenger
  prepares with probability $\zeta(\bar P)$ at each price $\bar P$ in the image of $P$;
\item bids: each prepared bidder bids its value.
\end{enumerate}
Entry after order $x$ is $e(x)=\rho+(1-\rho)\zeta(P(x))$. The tuple must satisfy six conditions.
\begin{enumerate}[label=(E\arabic*)]
\item \emph{Bayes consistency.} $\mu(x)=\pi f_H(x)/[\pi f_H(x)+(1-\pi)f_L(x)]$, where
  $f_j(x)=(1-\lambda)\eta(x)+\lambda\sigma_j(x)$ is the probability of order $x$ in class $j$.
\item \emph{Challenger beliefs.} At each price $\bar P$ in the image of $P$, the challenger's belief is
  $\bar\mu(\bar P)=\Pr(j=H\mid P(X)=\bar P)$, from the joint law of the class and the order.
\item \emph{Entry optimality.} $\zeta(\bar P)=1$ if $B(\bar\mu(\bar P))>c_H$, and $\zeta(\bar P)=0$ if
  $B(\bar\mu(\bar P))<c_H$. At equality the challenger tie rule applies.
\item \emph{Competitive prices.} $P(x)=\E[V_T\mid X=x]$ for every order $x$.
\item \emph{Investor optimality.} Against the fixed $P$ and $e$, the distribution $\sigma_j$ puts probability one
  on the maximizers of $U_j(q)=q\{\E[V_T\mid j,X=q]-P(q)\}-k|q|$. At indifference the investor tie rule applies.
\item \emph{Bids.} Truthful bids are weakly dominant in the second-price auction.
\end{enumerate}
\end{definition}

\paragraph{Tie rules.} The analysis states two rules for each indifference.
\begin{itemize}
\item \emph{Investor.} Rule IM (the baseline): an indifferent investor may use any mixture over its maximizers.
  This is the standard best-response condition. Rule I0: if abstention is a maximizer, the investor abstains.
\item \emph{High-cost challenger.} Rule C1 (the baseline, as in the Laplace manuscript): an indifferent
  challenger prepares. Rule CM: an indifferent challenger may prepare with any probability.
\end{itemize}
The zero-cost challenger prepares at every price under every rule. Its gross profit is an expected auction
rent, so it is never negative, and the model fixes preparation at zero gross profit.

\paragraph{Deviations and comparisons.} A unilateral deviation by the investor holds the candidate's $P$ and
$e$ fixed and changes only the order that reaches them. Lemma~\ref{lem:flow}(a) shows that every order has
positive probability in each class. A deviation therefore reaches a price in the image of $P$, and no
off-path belief enters any proof. A comparison across economies re-solves $P$, $\mu$ and $\zeta$ in each economy.
An \emph{outcome} is the tuple $(\sigma_H,\sigma_L,\mu,e,P)$; a unique outcome means a unique tuple.

%==========================================================================================
\section{Beliefs and prices}\label{sec:beliefs}

\begin{lemma}[flow law and beliefs; analytical]\label{lem:flow}
Let $0<\lambda<1$ and $0<\beta<1/2$, and fix any order distributions $\sigma_H,\sigma_L$.
\begin{enumerate}[label=(\alph*)]
\item Full support: $f_j(x)\ge(1-\lambda)\eta(x)>0$ for every order $x$ and class $j$.
\item Nonzero orders: for $x\in\{-1,1\}$, $b\le f_j(x)\le b+\lambda$, so $m(\pi)\le\mu(x)\le M(\pi)$.
\item The zero order: $b_0\le f_j(0)\le b_0+\lambda$. If $\beta\le1/3$, then $m(\pi)\le\mu(0)\le M(\pi)$.
  If $\beta>1/3$, the pair $\sigma_H=\delta_0$, $\sigma_L(0)=0$ gives $\mu(0)>M(\pi)$.
\item Full orders $\sigma_H=\delta_1$, $\sigma_L=\delta_{-1}$ give $(\mu(-1),\mu(0),\mu(1))=(m(\pi),\pi,M(\pi))$.
\end{enumerate}
\end{lemma}

\begin{proof}
(a) The liquidity trader fills the slot with probability $1-\lambda>0$ and places order $x$ with probability
$\eta(x)$. Each $\eta(x)$ is positive because $0<\beta<1/2$. The informed term $\lambda\sigma_j(x)$ is nonnegative.

(b) For $x=\pm1$, $(1-\lambda)\eta(x)=b$ and $0\le\lambda\sigma_j(x)\le\lambda$, so $b\le f_j(x)\le b+\lambda$.
Hence $f_H(x)/f_L(x)\in[b/(b+\lambda),(b+\lambda)/b]=[(1+u)^{-1},1+u]$. Bayes' rule in odds form is
\[
\frac{\mu(x)}{1-\mu(x)}=\frac{\pi}{1-\pi}\cdot\frac{f_H(x)}{f_L(x)}.
\]
The odds of $\mu(x)$ therefore lie between $\pi/[(1-\pi)(1+u)]$ and $\pi(1+u)/(1-\pi)$. These are the odds of
$m(\pi)$ and $M(\pi)$.

(c) For $x=0$, $(1-\lambda)\eta(0)=b_0$, so $b_0\le f_j(0)\le b_0+\lambda$, and the likelihood ratio lies in
$[(1+\lambda/b_0)^{-1},1+\lambda/b_0]$. The inequality $b_0\ge b$ is $1-2\beta\ge\beta$, that is $\beta\le1/3$.
Then $\lambda/b_0\le\lambda/b=u$, and the argument of (b) applies. If $\beta>1/3$, then $b_0<b$. The pair
$\sigma_H=\delta_0$, $\sigma_L(0)=0$ gives $f_H(0)=b_0+\lambda$ and $f_L(0)=b_0$, a ratio $1+\lambda/b_0>1+u$.
For example, $\lambda=6/17$, $\beta=2/5$ and $\pi=1/2$ give $\mu(0)=41/52>M=26/37$.

(d) Full orders give $f_H=(b,b_0,b+\lambda)$ and $f_L=(b+\lambda,b_0,b)$ on the orders $(-1,0,1)$.
The ratios are $(1+u)^{-1}$, $1$ and $1+u$.
\end{proof}

\paragraph{Where $\beta\le 1/3$ enters.} Only part (c) uses $\beta\le1/3$. Propositions 2--5 use parts (a), (b)
and (d), so they hold for every $\beta<1/2$. Part (c) extends the posterior bound to every order under every
strategy pair. Corollary~\ref{cor:ceiling} uses it.

\begin{lemma}[prices; analytical]\label{lem:price}
Suppose $\Delta>0$ and $g_H>g_L$. Fix a tuple that satisfies (E1)--(E4) and (E6), under either challenger tie rule.
\begin{enumerate}[label=(\alph*)]
\item Floor: $e(x)\ge\rho$ for every order $x$.
\item Price formula: $P(x)=t_0+e(x)[t_L-t_0+\Delta\mu(x)]$.
\item Sufficiency: if $P(x)=P(x')$, then $\mu(x)=\mu(x')$. The challenger's belief at the price $P(x)$ is $\mu(x)$.
\item Entry: $e(x)=\rho$ if $\mu(x)<\tau$ and $e(x)=1$ if $\mu(x)>\tau$. If $\mu(x)=\tau$, then $e(x)=1$ under C1
  and $e(x)\in[\rho,1]$ under CM.
\item Monotonicity: if $\mu(x)>\mu(x')$, then $e(x)\ge e(x')$ and $P(x)-P(x')\ge\rho\Delta[\mu(x)-\mu(x')]>0$.
\item Price algebra for a constructed tuple. Part (f) uses no equilibrium condition. Let
  $\mu:\{-1,0,1\}\to[0,1]$ and $e:\{-1,0,1\}\to[\rho,1]$ be any two functions such that $e(x)=e(x')$ when
  $\mu(x)=\mu(x')$, and $e(x)\ge e(x')$ when $\mu(x)>\mu(x')$. Let $t_L\ge t_0$, and define
  $P(x)=t_0+e(x)[t_L-t_0+\Delta\mu(x)]$. Then $P(x)=P(x')$ when $\mu(x)=\mu(x')$, and
  $P(x)-P(x')\ge\rho\Delta[\mu(x)-\mu(x')]>0$ when $\mu(x)>\mu(x')$. So $P(x)=P(x')$ exactly when $\mu(x)=\mu(x')$.
\end{enumerate}
\end{lemma}

\begin{proof}
(a) The zero-cost challenger prepares at every price, and it has probability $\rho$.

(b) Fix an order $x$. The challenger sees $P(x)$ and its cost. The cost is independent of the class, of $\theta$
and of $R$. Entry after $x$ therefore has probability $e(x)$, independent of $(j,\theta,R)$. With entry and
class $j$, expected proceeds are $t_j$; without entry they are $t_0$. So
$\E[V_T\mid j,X=x]=t_0+e(x)(t_j-t_0)$. Average over the class with weights $\mu(x)$ and $1-\mu(x)$, and use (E4).

(c) Let $P(x)=P(x')=\bar P$. The challenger cannot tell $x$ from $x'$, so $e(x)=e(x')=\bar e$, with $\bar e\ge\rho$
by (a). Part (b) gives $\bar e\Delta\mu(x)=\bar P-t_0-\bar e(t_L-t_0)=\bar e\Delta\mu(x')$. Since
$\bar e\Delta>0$, $\mu(x)=\mu(x')$. By (E2), $\bar\mu(\bar P)$ is an average of the posteriors of the orders with
price $\bar P$. These posteriors are equal, so $\bar\mu(\bar P)=\mu(x)$.

(d) By (c), the challenger's gross profit at $P(x)$ is $B(\mu(x))$. Since $g_H>g_L$, $B(\mu)>c_H$ exactly when
$\mu>\tau$, and $B(\mu)<c_H$ exactly when $\mu<\tau$. Apply (E3) and the tie rule.

(e) If $\mu(x')<\tau$, then $e(x')=\rho\le e(x)$. If $\mu(x')\ge\tau$, then $\mu(x)>\tau$ and $e(x)=1\ge e(x')$.
By (b),
\[
P(x)-P(x')=e(x)\Delta[\mu(x)-\mu(x')]+[e(x)-e(x')]\,[t_L-t_0+\Delta\mu(x')].
\]
The second term is nonnegative by (P1). The first is at least $\rho\Delta[\mu(x)-\mu(x')]$ by (a).

(f) If $\mu(x)=\mu(x')$, then $e(x)=e(x')$, and the two prices have the same formula. If $\mu(x)>\mu(x')$, the
identity in (e) holds for the defined $P$:
\[
P(x)-P(x')=e(x)\Delta[\mu(x)-\mu(x')]+[e(x)-e(x')]\,[t_L-t_0+\Delta\mu(x')].
\]
Here $e(x)\ge\rho$, $\Delta>0$, $e(x)-e(x')\ge0$, $t_L-t_0\ge0$ and $\Delta\mu(x')\ge0$. So the difference is at
least $\rho\Delta[\mu(x)-\mu(x')]>0$. The argument uses only these inequalities, not (E1)--(E4).
\end{proof}

Parts (a)--(e) apply to a tuple that is already known to satisfy (E1)--(E4). A proof that \emph{constructs} a
tuple and then verifies (E2)--(E4) uses part (f) instead. Each verification below follows this order:
the price formula defines $P$; part (f) shows which orders share a price; (E2) then gives the challenger's
belief at each price; and (E3) and (E4) follow by direct computation.

%==========================================================================================
\section{Informed trading}\label{sec:trading}

\begin{lemma}[informed trading; analytical]\label{lem:trade}
Assume $\Delta\ge0$, as (P1) gives. Fix a tuple of the form in Definition~\ref{def:eq} with these properties:
\begin{itemize}
\item its posterior $\mu$ satisfies Bayes consistency (E1) for the order distributions $\sigma_H,\sigma_L$;
\item its price $P$ satisfies competitive pricing (E4);
\item its entry $e(x)=\rho+(1-\rho)\zeta(P(x))$ lies in $[\rho,1]$ after every order. This floor holds for every
  tuple of Definition~\ref{def:eq}, because the zero-cost challenger prepares at every price.
\end{itemize}
The payoff formulas use (E4) and the entry $e$ only. Part (c) also uses (E1) and the floor.
Conditional on its arrival, class $j$ earns these expected payoffs:
\[
\begin{aligned}
U_H(1)&=e(1)\Delta[1-\mu(1)]-k,&\quad U_H(0)&=0,&\quad U_H(-1)&=-e(-1)\Delta[1-\mu(-1)]-k,\\
U_L(-1)&=e(-1)\Delta\mu(-1)-k,&\quad U_L(0)&=0,&\quad U_L(1)&=-e(1)\Delta\mu(1)-k.
\end{aligned}
\]
Write $K_H=e(1)\Delta[1-\mu(1)]$ and $K_L=e(-1)\Delta\mu(-1)$ for the gross returns of the two correct orders.
\begin{enumerate}[label=(\alph*)]
\item A wrong-sign order earns at most $-k<0$, so abstention strictly beats it.
\item Class $H$ buys with probability one if $K_H>k$ and abstains if $K_H<k$. If $K_H=k$, the purchase and
  abstention are its two maximizers. Class $L$ sells with probability one if $K_L>k$ and abstains if $K_L<k$.
  If $K_L=k$, the sale and abstention are its two maximizers.
\item Every candidate satisfies $\rho\Delta[1-M(\pi)]\le K_H\le\Delta$ and $\rho\Delta m(\pi)\le K_L\le\Delta$.
  Hence $\rho\kappa(\pi)\Delta\le K_j\le\Delta$.
\item $K_H$ depends on the candidate only through $(\mu(1),e(1))$, and $K_L$ only through $(\mu(-1),e(-1))$.
  Entry after the zero order enters no payoff.
\end{enumerate}
\end{lemma}

\begin{proof}
The order $q$ reaches the price $P(q)$. As in Lemma~\ref{lem:price}(b), $\E[V_T\mid j,X=q]=t_0+e(q)(t_j-t_0)$.
Condition (E4) gives $P(q)=t_0+e(q)[t_L-t_0+\Delta\mu(q)]$; this step uses only competitive prices and
needs no condition on $\Delta$. Subtract:
\[
\E[V_T\mid H,X=q]-P(q)=e(q)\Delta[1-\mu(q)],\qquad P(q)-\E[V_T\mid L,X=q]=e(q)\Delta\mu(q).
\]
Multiply by $q$ and subtract $k|q|$. This gives the six payoffs. The terms $t_0$ and $e(q)(t_L-t_0)$ cancel.

(a) $e(q)\ge0$, $\Delta\ge0$ and $\mu(q)\in[0,1]$, so each wrong-sign payoff is at most $-k$.

(b) By (a), each class compares its correct order with abstention, and $U_j(\pm1)-U_j(0)=K_j-k$.

(c) The entry satisfies $\rho\le e\le1$ by the premise. The posterior satisfies (E1), so
Lemma~\ref{lem:flow}(b) gives $1-\mu(1)\ge1-M(\pi)$ and $\mu(-1)\ge m(\pi)$. Also $1-\mu(1)\le1$ and $\mu(-1)\le1$.
Multiply by $e\Delta$, with $\Delta\ge0$.

(d) Read off the two formulas.
\end{proof}

The payoffs are conditional on the investor's arrival. The probability $\lambda$ multiplies both sides of every
comparison, so the comparison of $K_j$ with $k$ is the single decision of each class.

\begin{corollary}[the ceiling under every profile; analytical]\label{cor:ceiling}
Let $\beta\le1/3$ and $g_H\ge g_L$. If $B(M(\pi))<c_H$, then no candidate profile, in or out of equilibrium,
has high-cost preparation after any order. Here a candidate profile is any pair $(\sigma_H,\sigma_L)$ with the
beliefs that (E1) and (E2) give and the preparation rule that (E3) gives.
\end{corollary}

\begin{proof}
Lemma~\ref{lem:flow}(b)--(c) give $\mu(x)\le M(\pi)$ at every order under every strategy pair. The challenger's
belief at a price is an average of such posteriors, by (E2). Since $B$ is nondecreasing,
$B(\bar\mu)\le B(M(\pi))<c_H$ at every price, and (E3) gives $\zeta=0$ at every price.
\end{proof}

%==========================================================================================
\section{Proposition 2: competition creates competition}\label{sec:prop2}

\setcounter{proposition}{1}
\begin{proposition}[competition creates competition; analytical]\label{prop:2}
Let the classes satisfy the restrictions of Section 1.1. Let $F_1$ be a strengthening of $F_0$, and $F_2$ a
strengthening of $F_1$. Fix $0<\pi<1$, $0<\rho<1$, $k>0$, $0<\lambda<1$ and $0<\beta<1/2$. Write $B_i=B_{F_i}$,
$\Delta_i=\Delta(F_i)$, $M=M(\pi)$, $m=m(\pi)$ and $\kappa=\kappa(\pi)$. Suppose
\begin{align*}
&\text{(H1)}\quad B_0(\pi)<c_H<B_1(M),\\
&\text{(H2)}\quad \Delta_0<k<\rho\kappa\Delta_1,\\
&\text{(H3)}\quad B_2(M)<c_H.
\end{align*}
Then each economy has a unique equilibrium outcome. The outcome is the same under all four pairs of tie rules.
\begin{enumerate}[label=(\roman*)]
\item \emph{Weak incumbent $F_0$.} No informed trade: $\sigma_H=\sigma_L=\delta_0$. Every posterior equals $\pi$,
  entry equals $\rho$ after every order, and every order has the price $t_0+\rho[t_L-t_0+\Delta_0\pi]$.
  The price carries no information.
\item \emph{Strong incumbent $F_1$.} Full two-sided trade: $\sigma_H=\delta_1$ and $\sigma_L=\delta_{-1}$. The
  posteriors after the orders $(-1,0,1)$ are $(m,\pi,M)$. Only a purchase recruits the high-cost challenger,
  so $e=(\rho,\rho,1)$. The prices are $P(x)=t_0+e(x)[t_L-t_0+\Delta_1\mu(x)]$, with $P(-1)<P(0)<P(1)$. Entry is
  \[
  \Ent_1=\rho+(1-\rho)(b+\pi\lambda)>\rho .
  \]
  Entry given class $H$ is $e_H=\rho+(1-\rho)(b+\lambda)$, and given class $L$ it is $e_L=\rho+(1-\rho)b$.
  High-class ownership, the probability that a class-$H$ challenger acquires the target, is
  $\Own_{H,1}=\pi e_HW_H(F_1)$. It exceeds its weak-incumbent level $\Own_{H,0}=\pi\rho W_H(F_0)$ if and only if
  \[
  [\rho+(1-\rho)(b+\lambda)]\,W_H(F_1)>\rho\,W_H(F_0).
  \]
\item \emph{Very strong incumbent $F_2$.} Full two-sided trade again, with posteriors $(m,\pi,M)$. Entry equals
  $\rho$ after every order. The three prices $t_0+\rho[t_L-t_0+\Delta_2\mu(x)]$ differ, so the price stays
  informative while entry returns to $\rho$.
\end{enumerate}
\end{proposition}

In the point-mass benchmark $h$ exceeds the incumbent's support at every strength, so $W_H=1$ and ownership rises.
The Laplace manuscript states the sufficient condition $(1+\rho)W_H(F_1)\ge2\rho W_H(F_0)>0$. It implies the
condition of part (ii) when $b+\lambda>1/2$, because then $e_H>(1+\rho)/2$. At calibration B, $b+\lambda=0.531737$.

\begin{proof}
\emph{Step 0: implied inequalities.} By (P2), $B_1(\pi)\le B_0(\pi)<c_H$, so (H1) gives
$B_1(\pi)<c_H<B_1(M)$. Since $M>\pi$, $B_1(M)-B_1(\pi)=(M-\pi)[g_H(F_1)-g_L(F_1)]>0$, so $g_H(F_1)>g_L(F_1)$ and
$\pi<\tau_1<M$, where $\tau_1$ is the threshold at $F_1$. (H2) gives $\Delta_1>k/(\rho\kappa)>0$. By (P3),
$\Delta_2\ge\Delta_1$, so $\rho\kappa\Delta_2>k$. By (P1), $g_H(F_2)\ge g_L(F_2)$, so
$B_2(m)\le B_2(\pi)\le B_2(M)<c_H$.

\emph{Step 1: the weak incumbent, every candidate.} Fix any candidate tuple at $F_0$. By Lemma~\ref{lem:trade},
a correct order earns $K_j-k\le\Delta_0-k<0$, and a wrong-sign order earns at most $-k$. Abstention is the unique
maximizer of each class, so every equilibrium has $\sigma_H=\sigma_L=\delta_0$, under both investor rules.
Then $f_H(x)=f_L(x)=(1-\lambda)\eta(x)+\lambda\one\{x=0\}>0$: the informed investor's abstention adds $\lambda$ to
the zero-order mass in both classes. So (E1) gives $\mu(x)=\pi$ after every order. All orders have the
same posterior, so the challenger's belief at every price is $\pi$ by (E2). Since $B_0(\pi)<c_H$, (E3) gives
$\zeta=0$ at every price, with no tie. Entry is $\rho$ after every order. As in Lemma~\ref{lem:price}(b), (E4)
gives $P(x)=t_0+\rho[\pi t_H+(1-\pi)t_L-t_0]=t_0+\rho[t_L-t_0+\Delta_0\pi]$ for every $x$.

\emph{Verification of the tuple in part (i).} Define the tuple: $\sigma_H=\sigma_L=\delta_0$, $\mu\equiv\pi$,
$\zeta=0$ at the single price $\bar P=t_0+\rho[t_L-t_0+\Delta_0\pi]$, $P(x)=\bar P$ for every $x$, and truthful bids.
(E1) holds because $f_H=f_L$. (E2): every order has the price $\bar P$, so the challenger's belief at $\bar P$ is
the prior $\pi$. (E3) holds strictly, since $B_0(\pi)<c_H$. (E4): entry is $\rho$ after every order and is
independent of $(j,\theta,R)$, so $\E[V_T\mid X=x]=t_0+\rho[\pi t_H+(1-\pi)t_L-t_0]=\bar P$. (E5): the tuple
satisfies the premises of Lemma~\ref{lem:trade}, and the bound of Step 1 makes abstention strictly optimal.
(E6) holds. Step 1 shows that every equilibrium is this tuple.

\emph{Step 2: the strong incumbent, orders in every candidate.} Fix any candidate tuple at $F_1$. By
Lemma~\ref{lem:trade}(c), $K_H\ge\rho\Delta_1(1-M)\ge\rho\kappa\Delta_1>k$ and $K_L\ge\rho\Delta_1m\ge\rho\kappa\Delta_1>k$.
By Lemma~\ref{lem:trade}(a)--(b), the purchase is the unique maximizer of class $H$, and the sale is the unique
maximizer of class $L$. Every equilibrium has $\sigma_H=\delta_1$ and $\sigma_L=\delta_{-1}$, under both
investor rules. The bounds use Lemma~\ref{lem:flow}(b) only, so they hold for every $\beta<1/2$.

\emph{Step 3: the strong incumbent, the tuple.} \emph{Necessity.} Fix any equilibrium at $F_1$. Step 2 gives
full orders, and Lemma~\ref{lem:flow}(d) gives the posteriors $(m,\pi,M)$. Step 0 gives $\Delta_1>0$ and
$g_H(F_1)>g_L(F_1)$, so Lemma~\ref{lem:price} applies to this equilibrium. Since $B_1(M)>c_H$, $B_1(\pi)<c_H$ and
$B_1(m)\le B_1(\pi)$, Lemma~\ref{lem:price}(d) gives $e(1)=1$ and $e(0)=e(-1)=\rho$, with no tie.
Lemma~\ref{lem:price}(b) gives the prices. So every equilibrium is the tuple of part (ii).

\emph{Verification.} Define the tuple of part (ii): $\sigma_H=\delta_1$, $\sigma_L=\delta_{-1}$,
$\mu=(m,\pi,M)$ on the orders $(-1,0,1)$, $e=(\rho,\rho,1)$, $P(x)=t_0+e(x)[t_L-t_0+\Delta_1\mu(x)]$, $\zeta=1$ at
$P(1)$ and $\zeta=0$ at $P(0)$ and $P(-1)$, and truthful bids. The verification does not use Lemma~\ref{lem:price}(a)--(e).
\begin{itemize}
\item (E1) holds by Lemma~\ref{lem:flow}(d).
\item (E2): $e$ is nondecreasing in $\mu$, $m<\pi<M$, $\Delta_1>0$ and $t_L\ge t_0$ by (P1). So
  Lemma~\ref{lem:price}(f) gives $P(-1)<P(0)<P(1)$. Each price comes from one order, and the challenger's belief at
  $P(x)$ is $\mu(x)$.
\item (E3): $B_1(M)>c_H$ gives $\zeta=1$ at $P(1)$. $B_1(m)\le B_1(\pi)<c_H$ gives $\zeta=0$ at $P(0)$ and
  $P(-1)$. Each comparison is strict, and $\zeta$ gives the stated $e$.
\item (E4): after the order $x$, entry has probability $e(x)$, independent of $(j,\theta,R)$, because the cost is
  independent of them. So $\E[V_T\mid X=x]=t_0+e(x)[\mu(x)t_H+(1-\mu(x))t_L-t_0]=P(x)$.
\item (E5): the tuple satisfies the premises of Lemma~\ref{lem:trade}, and the bounds of Step 2 make the full
  orders strictly optimal.
\item (E6) holds.
\end{itemize}

\emph{Step 4: entry and ownership at $F_1$.} Class $H$ places the order $1$ with probability $b+\lambda$: a
liquidity purchase or an informed purchase. Class $L$ places it with probability $b$. So
$e_H=\rho+(1-\rho)(b+\lambda)$, $e_L=\rho+(1-\rho)b$ and $\Ent_1=\pi e_H+(1-\pi)e_L=\rho+(1-\rho)(b+\pi\lambda)$.
Since $b>0$, $\Ent_1>\rho$. Given class $H$ and preparation, the challenger wins when $\theta>R$; ties have
probability zero because $F$ is continuous. Preparation depends on the order and the cost, and both are
independent of $(\theta,R)$ given the class. So a class-$H$ challenger acquires the target with probability
$e_HW_H(F_1)$, and $\Own_{H,1}=\pi e_HW_H(F_1)$. At $F_0$ entry is $\rho$ in each class, so
$\Own_{H,0}=\pi\rho W_H(F_0)$. Divide the comparison by $\pi>0$.

\emph{Step 5: the very strong incumbent.} Step 0 gives $\rho\kappa\Delta_2>k$. The argument of Step 2 with
$\Delta_2$ in place of $\Delta_1$ gives full orders in every equilibrium, and Lemma~\ref{lem:flow}(d) gives the
posteriors $(m,\pi,M)$. Step 0 gives $B_2(\mu)\le B_2(M)<c_H$ at each of the three posteriors, so $\zeta=0$
at every price, with no tie. Entry is $\rho$ after every order. \emph{Verification.} Define the tuple of part
(iii): full orders, $\mu=(m,\pi,M)$, $e\equiv\rho$, $P(x)=t_0+\rho[t_L-t_0+\Delta_2\mu(x)]$, $\zeta=0$ at every
price, and truthful bids. Lemma~\ref{lem:price}(f) with a constant $e$, $\Delta_2>0$ and $m<\pi<M$ gives three
distinct prices, so the belief at $P(x)$ is $\mu(x)$. Then (E1), (E3) (strictly, by Step 0), (E4), (E5) and (E6)
follow as in Step 3, with the bounds of Step 2 at $\Delta_2$.

\emph{Step 6: tie rules.} In each economy every investor comparison is strict (Steps 1, 2 and 5). Every
equilibrium posterior lies strictly on one side of the threshold (Steps 1, 3 and 5). No tie rule is ever used,
so the four pairs of rules give the same unique outcome.
\end{proof}

\paragraph{The hypotheses are complete, and each one is used.} The proof uses each hypothesis once.
\begin{itemize}
\item $B_0(\pi)<c_H$ gives entry $\rho$ at $F_0$. Through (P2), it also stops recruitment after the zero order at $F_1$.
\item $c_H<B_1(M)$ gives recruitment after a purchase at $F_1$.
\item $\Delta_0<k$ defeats every informed order at $F_0$, against every candidate.
\item $k<\rho\kappa\Delta_1$ forces full orders at $F_1$, and through (P3) at $F_2$, against every candidate.
\item $B_2(M)<c_H$ stops recruitment at $F_2$.
\end{itemize}
Nothing else enters the proof; in particular Lemma~\ref{lem:flow}(c), the only use of $\beta\le1/3$, does not.
The hypotheses imply $B_1(\pi)<c_H$, $\rho\kappa\Delta_2>k$ and $B_2(m)\le B_2(\pi)\le B_2(M)$.

A list that replaces $B_0(\pi)<c_H$ by the weaker $B_1(\pi)<c_H$ is not sufficient. Take calibration B and the
benchmark values, but set $c_H=47/10$ and $(r_0,r_1,r_2)=(6/5,3,8)$. Then $k-\Delta_0=1/300$,
$c_H-B_1(1/2)=49/120$, $B_1(M)-c_H=1.517155$, $\rho\kappa\Delta_1-k=0.024824$ and $c_H-B_2(M)=0.312465$ are positive,
but $c_H-B_0(1/2)=-5/48$. At $r_0$ the orders are zero by Step 1, the posterior is $1/2$, and
$B_0(1/2)=1153/240>47/10$. Both cost types prepare, so entry is $1$ at $r_0$. Entry then falls to
$\Ent_1=0.522757$ at $r_1$.

On the benchmark path each other hypothesis also guards a conclusion. At $r_0=31/20$ the left side of (H2)
fails, since $\Delta(31/20)=121/1240>k$, and Proposition~\ref{prop:3} gives two equilibria there. At $r_1=2$ the
right side of (H2) fails, since $\rho\kappa\Delta(2)-k=-0.003191$, and Proposition~\ref{prop:3} gives a mixed sale
there. If $F_1=F_2$, (H3) or the right side of (H1) fails, and part (iii) or part (ii) fails with it.

\paragraph{The open set.} For uniform incumbents, (H1)--(H3) are strict inequalities between continuous
functions of $(r_0,r_1,r_2,c_H,k,\rho,\pi,\lambda,\beta)$. They define an open set. Proposition~\ref{prop:region}
gives its section through the benchmark.

\paragraph{Calibration B.} Table~\ref{tab:prop2} lists the three equilibrium tuples at calibration B and prior
$1/2$. The values come from the script \texttt{p1\_numbers.py} in outward interval arithmetic.

\begin{table}[h]
\centering
\caption{The equilibrium tuples of Proposition~\ref{prop:2} at calibration B, prior $1/2$.}\label{tab:prop2}
\small
\setlength{\tabcolsep}{4pt}
\begin{tabular}{@{}lccc@{}}
\toprule
 & $r_0=1.2$ & $r_1=3$ & $r_2=3.6$\\
\midrule
Orders $(\sigma_H,\sigma_L)$ & $(\delta_0,\delta_0)$ & $(\delta_1,\delta_{-1})$ & $(\delta_1,\delta_{-1})$\\
Posteriors after $(-1,0,1)$ & $(0.5,0.5,0.5)$ & $(0.268941,0.5,0.731059)$ & $(0.268941,0.5,0.731059)$\\
Preparation, high cost & $(0,0,0)$ & $(0,0,1)$ & $(0,0,0)$\\
Prices after $(-1,0,1)$ & $0.392708$ each & $(0.576074,0.614583,1.362372)$ & $(0.610002,0.664236,0.718471)$\\
Entry & $0.250000$ & $0.522757$ & $0.250000$\\
High-class ownership & $0.125000$ & $0.324401$ & $0.125000$\\
\bottomrule
\end{tabular}
\end{table}

\section{Proposition 3: the equilibrium set along the benchmark path}\label{sec:prop3}

Throughout this section the auction is the benchmark: point masses $h$ and $\ell$, $R\sim U[0,r]$, and the
benchmark values of $(h,\ell,p,\rho,c_H,k)$. The path is $r\in[r_0,r_2]=[6/5,18/5]$. The slot is any pair
$(\lambda,\beta)$ with $0<\lambda<1$, $0<\beta<1/2$ and $\tau(r_1)<M<\tau(r_2)$, the band of
Proposition~\ref{prop:region}. Section 5.3 takes the prior $\pi=1/2$; Sections 5.2 and 5.4 allow any prior.

\subsection{The path and its thresholds}

Two functions organize the path. The first is
\[
J(r)=\Delta(r)\,[1-\tau(r)]=\frac{(r-\ell)^2\,[2r(h-c_H)-r^2-p^2]}{2r\,(2rh-r^2-\ell^2)},
\]
the gross return of a purchase that reaches the posterior $\tau(r)$ and recruits. The second is
$\fr(d)=\ell+d+\sqrt{d^2+2\ell d}$, the root above $\ell$ of $\Delta(r)=d$. Since
$\Delta'(r)=\tfrac12-\ell^2/(2r^2)>0$ for $r>\ell$, $\Delta(r)\le d$ holds exactly when $r\le\fr(d)$.

\begin{lemma}[the path; computer-assisted]\label{lem:path}
\begin{enumerate}[label=(\alph*)]
\item $\tau$ is strictly increasing on $[r_0,r_2]$, with $\tau(r_0)=195/308>1/2$ and $\tau(r_2)=4245/5804<3/4$.
\item $J$ is strictly increasing on $[r_0,r_2]$, with $J(r_0)=113/18480<k$ and $\rho J(r_2)=263471/4178880>k$.
\item $\Delta(r)[1-\tau(r)]^2$ is strictly increasing on $[r_0,2]$.
\item $\Delta(r)\tau(r)$ is strictly increasing on $[r_0,r_2]$.
\end{enumerate}
\end{lemma}

\begin{proof}
(a) Write $\tau=N_\tau/D_\tau$ with $N_\tau=2rc_H-\ell^2+p^2$ and $D_\tau=2rh-r^2-\ell^2$. The derivative has the
sign of $N_\tau'D_\tau-N_\tau D_\tau'=2c_Hr^2-2(\ell^2-p^2)r+2h(\ell^2-p^2)-2c_H\ell^2$, which at the benchmark is
$12r^2-\tfrac32r+3$. Its discriminant is negative, so it is positive for every $r$. Also $D_\tau>0$ on the path.

(b) $J'$ has the sign of
\[
Q_J(r)=4r^6-160r^5+743r^4-12r^3-614r^2+40r-1 .
\]
The script \texttt{p1\_certificates.py} computes $Q_J$ from the primitives with exact rational arithmetic. A Sturm
sequence gives no real root of $Q_J$ in $(r_0,r_2]$, and $Q_J(r_0)=4634399/15625>0$. The two values follow
from the closed form.

(c) The derivative of $\Delta(1-\tau)^2$ has the sign of the degree-10 polynomial
\[
16r^{10}-1280r^9+29256r^8-229344r^7+562917r^6-32808r^5-374607r^4+48456r^3-2685r^2+80r-1 .
\]
A Sturm sequence gives no real root in $(r_0,2]$, and its value at $r_0$ is positive (script
\texttt{p1\_certificates.py}).

(d) $\Delta$ and $\tau$ are positive and strictly increasing on the path, by (a) and $\Delta'>0$.

Parts (a) and (d) are proved by hand. Parts (b) and (c) rest on the exact Sturm counts.
\end{proof}

\begin{definition}[thresholds]\label{def:thresholds}
At the prior $1/2$, with $M=(1+u)/(2+u)$ and $m=1-M$:
\begin{itemize}
\item $r_A$: the root of $J(r)=k$ in $(r_0,r_2)$. A purchase that just reaches $\tau$ and recruits earns $k$.
\item $r_N=\fr(2k/\rho)$. Against no informed trade, a unit order earns $\rho\Delta/2$; this equals $k$ at $r_N$.
\item $r_F=\fr\big(k/(1-M)\big)$. A full purchase with recruitment earns $\Delta(1-M)$; this equals $k$ at $r_F$.
\item $r_J$: the root of $\rho J(r)=k$ in $(r_0,r_2)$. A purchase that just reaches $\tau$ without recruitment
  earns $\rho J$; this equals $k$ at $r_J$.
\item $r_S=\fr\big(k/(\rho m)\big)$. A full sale earns $\rho\Delta m$; this equals $k$ at $r_S$.
\item $r_C$: the root of $\tau(r)=M$, that is of $B_r(M)=c_H$. In closed form,
  \[
  r_C=\frac{Mh-c_H+\sqrt{(Mh-c_H)^2+M[(1-M)\ell^2-p^2]}}{M}.
  \]
\end{itemize}
\end{definition}

Lemma~\ref{lem:path}(a)--(b) give the uniqueness of $r_A$, $r_J$ and $r_C$. The thresholds $r_A$, $r_N$ and $r_J$
do not depend on the slot. At the benchmark $r_N=(29+6\sqrt6)/25$.

\begin{lemma}[the order of the thresholds; computer-assisted]\label{lem:order}
For every slot in the band, $r_0<r_A<r_F<r_N<r_J<r_S<r_1<r_C<r_2$.
\end{lemma}

\begin{proof}
Each inequality reduces to a fact already stated.
\begin{itemize}
\item $r_0<r_A<r_1$: $J(r_0)<k<\rho J(r_1)<J(r_1)$, with $\rho J(r_1)=59/1200$, and $J$ is increasing.
\item $r_A<r_F$: $\Delta(r_A)<k/(1-M)$ is $(1-M)\Delta(r_A)<k=(1-\tau(r_A))\Delta(r_A)$, that is $M>\tau(r_A)$.
  This holds because $M>\tau(r_1)>\tau(r_A)$, since $r_A<r_1$.
\item $r_F<r_N$: $k/(1-M)<2k/\rho$ is $M<1-\rho/2=7/8$. This holds because $M<\tau(r_2)<3/4$.
\item $r_N<r_J$: $\rho J(r_N)=\rho\Delta(r_N)[1-\tau(r_N)]=2k[1-\tau(r_N)]<k$, because $\tau(r_N)>1/2$.
\item $r_J<r_S$: $\Delta(r_J)<k/(\rho m)$ is $\rho(1-M)\Delta(r_J)<k=\rho(1-\tau(r_J))\Delta(r_J)$, that is
  $M>\tau(r_J)$. This holds because $\rho J(r_1)=59/1200>k$ gives $r_J<r_1$, and then $M>\tau(r_1)>\tau(r_J)$.
\item $r_S<r_1$: $\rho(1-M)\Delta(r_1)>k$ is $(1-M)/6>1/50$, that is $M<22/25$. This holds because $M<3/4$.
\item $r_1<r_C<r_2$: this is $\tau(r_1)<M<\tau(r_2)$, the band, since $\tau$ is increasing.\qedhere
\end{itemize}
\end{proof}

Table~\ref{tab:thresholds} gives the thresholds at calibration B.

\begin{table}[h]
\centering
\caption{Thresholds of the benchmark path at calibration B, prior $1/2$ (\lab{computer-assisted}).}\label{tab:thresholds}
\small
\begin{tabular}{@{}llll@{}}
\toprule
Threshold & Definition & Value at B & Depends on the slot\\
\midrule
$r_A$ & $J(r)=k$ & $1.393958$ & no\\
$r_F$ & $\Delta(r)(1-M)=k$ & $1.467127$ & yes\\
$r_N$ & $\rho\Delta(r)/2=k$ & $1.747878$ & no\\
$r_J$ & $\rho J(r)=k$ & $1.966097$ & no\\
$r_S$ & $\rho\Delta(r)m=k$ & $2.124148$ & yes\\
$r_C$ & $\tau(r)=M$ & $3.592659$ & yes\\
\bottomrule
\end{tabular}
\end{table}

\subsection{The two investor problems}

Fix a strength $r$ on the path and a prior $\pi$ with $\pi<\tau(r)$. By Lemma~\ref{lem:trade}(a), no equilibrium
puts weight on a wrong-sign order. Write $s=\lambda\sigma_H(1)$ and $t=\lambda\sigma_L(-1)$ for the probabilities
of an informed purchase and an informed sale. On the orders $(-1,0,1)$,
\[
f_H=(b,\;b_0+\lambda-s,\;b+s),\qquad f_L=(b+t,\;b_0+\lambda-t,\;b).
\]
Bayes' rule gives
\[
\mu_+(s)=\mu(1)=\frac{\pi(b+s)}{b+\pi s},\qquad 1-\mu_+(s)=\frac{(1-\pi)b}{b+\pi s},\qquad
\mu_-(t)=\mu(-1)=\frac{\pi b}{b+(1-\pi)t},
\]
and the posterior $\mu_0(s,t)=\mu(0)$ has the odds $\frac{\pi}{1-\pi}\cdot\frac{b_0+\lambda-s}{b_0+\lambda-t}$.
On $[0,\lambda]$, $\mu_+$ is continuous and strictly increasing from $\pi$ to $M(\pi)$, and $\mu_-$ is continuous
and strictly decreasing from $\pi$ to $m(\pi)$. When $\pi<\tau<M(\pi)$, the rate
\[
s^*=\frac{b(\tau-\pi)}{\pi(1-\tau)}\in(0,\lambda)
\]
solves $\mu_+(s^*)=\tau$. In Lemmas~\ref{lem:H} and~\ref{lem:L}, ``(E3) after the order $x$'' means that $e(x)$ is
optimal for the belief $\mu(x)$: $e(x)=\rho$ if $\mu(x)<\tau$, $e(x)=1$ if $\mu(x)>\tau$, and the tie rule at
equality. In an equilibrium this is Lemma~\ref{lem:price}(d). In a constructed tuple, Lemma~\ref{lem:product} shows
that the challenger's belief at the price of $x$ is $\mu(x)$. On the path, $\Delta(r)\ge\Delta(r_0)=1/60>0$ and
$g_H-g_L\ge1451/180>0$, so Lemma~\ref{lem:price} applies at every strength.

\begin{lemma}[the high class; analytical]\label{lem:H}
Fix $r$ and $\pi$ with $\pi<\tau(r)$. The pair $(s,e(1))$ satisfies (E1), (E3) after the order $1$ and (E5) for
class $H$ exactly when it is one of these five pairs.
\begin{enumerate}[label=(\Alph*)]
\item[(N)] $s=0$ and $e(1)=\rho$, with $\rho\Delta(1-\pi)\le k$.
\item[(Q)] $0<s<\min\{s^*,\lambda\}$, $e(1)=\rho$ and $\rho\Delta[1-\mu_+(s)]=k$. Only under IM.
\item[(R)] $s^*<s<\lambda$, $e(1)=1$ and $\Delta[1-\mu_+(s)]=k$. Only under IM.
\item[(T)] $s=s^*\in(0,\lambda)$ and $e(1)=e^*:=k/J(r)$. Under C1 this requires $e^*=1$, that is $J(r)=k$.
  Under CM it requires $\rho\le e^*\le1$. Only under IM.
\item[(F)] $s=\lambda$. Here $e(1)=1$ if $M(\pi)>\tau$ and $e(1)=\rho$ if $M(\pi)<\tau$. If $M(\pi)=\tau$,
  then $e(1)=1$ under C1, and $e(1)\in[\rho,1]$ under CM. The return must satisfy $e(1)\Delta[1-M(\pi)]\ge k$
  under IM, and $>k$ under I0.
\end{enumerate}
Each equation in (Q) and (R) has at most one solution.
\end{lemma}

\begin{proof}
By Lemma~\ref{lem:trade}(b), (E5) for class $H$ requires $K_H\le k$ at $s=0$; $K_H\ge k$ at $s=\lambda$ under IM,
and $K_H>k$ under I0; and $K_H=k$ at $0<s<\lambda$. An interior $s$ is a mixture at indifference, so it requires
IM; under I0 indifference forces $s=0$. Here $K_H=e(1)\Delta[1-\mu_+(s)]$, and $e(1)$ follows from
the entry rule at the posterior $\mu_+(s)$.
\begin{itemize}
\item At $s=0$, $\mu_+(0)=\pi<\tau$, so $e(1)=\rho$. This is (N).
\item At $0<s<\lambda$ with $\mu_+(s)<\tau$: $e(1)=\rho$, and $K_H=k$ is the equation of (Q).
\item At $0<s<\lambda$ with $\mu_+(s)>\tau$: $e(1)=1$, and $K_H=k$ is the equation of (R).
\item At $0<s<\lambda$ with $\mu_+(s)=\tau$: then $s=s^*$, and $K_H=e(1)\Delta(1-\tau)=e(1)J(r)=k$ gives
  $e(1)=e^*$. Under C1, $e(1)=1$. Under CM, $e(1)$ may be any value in $[\rho,1]$. This is (T).
\item At $s=\lambda$, $\mu_+(\lambda)=M(\pi)$. The entry rule gives $e(1)$, and the condition on
  $K_H$ is that of (F).
\end{itemize}
Since $\mu_+$ is strictly increasing, $1-\mu_+$ is strictly decreasing, so each equation has at most one root.
\end{proof}

\begin{lemma}[the low class; analytical]\label{lem:L}
Fix $r$ and $\pi$ with $\pi<\tau(r)$. Then $\mu_-(t)\le\pi<\tau$ for every $t$, so $e(-1)=\rho$ in every
equilibrium. The rate $t$ satisfies (E1), (E3) after the order $-1$ and (E5) for class $L$ if and only if it is
one of these:
\begin{enumerate}
\item[(N)] $t=0$, with $\rho\Delta\pi\le k$;
\item[(Q)] $0<t<\lambda$ with $\rho\Delta\mu_-(t)=k$; only under IM;
\item[(F)] $t=\lambda$, with $\rho\Delta m(\pi)\ge k$ under IM, and $>k$ under I0.
\end{enumerate}
Under IM exactly one rate qualifies at each $r$:
\[
\bar t(r)=\begin{cases}0, & \rho\Delta\pi\le k,\\[2pt]
\dfrac{b\,(\pi\rho\Delta/k-1)}{1-\pi}, & \rho\Delta m(\pi)<k<\rho\Delta\pi,\\[6pt]
\lambda, & \rho\Delta m(\pi)\ge k.\end{cases}
\]
Under I0 no rate qualifies when $\rho\Delta m(\pi)\le k<\rho\Delta\pi$.
\end{lemma}

\begin{proof}
Since $\mu_-$ is decreasing, $\mu_-(t)\le\mu_-(0)=\pi<\tau$, so the entry rule gives $e(-1)=\rho$. Then
$K_L(t)=\rho\Delta\mu_-(t)$, a continuous and strictly decreasing function of $t$, from $\rho\Delta\pi$ at $t=0$ to
$\rho\Delta m(\pi)$ at $t=\lambda$. By Lemma~\ref{lem:trade}(b), (E5) for class $L$ requires $K_L(0)\le k$ at
$t=0$; $K_L(\lambda)\ge k$ at $t=\lambda$ under IM, and $>k$ under I0; and $K_L(t)=k$ at $0<t<\lambda$, under IM
only. If $\rho\Delta\pi\le k$, then $K_L(t)<k$ for every $t>0$, so only $t=0$ qualifies. If $\rho\Delta m(\pi)\ge k$,
then $K_L(t)>k$ for every $t<\lambda$, so only $t=\lambda$ qualifies. Otherwise $\rho\Delta m(\pi)<k<\rho\Delta\pi$,
and the intermediate value theorem gives exactly one interior root. The equation
$\pi b/[b+(1-\pi)t]=k/(\rho\Delta)$ gives the formula. Under I0 with $\rho\Delta m(\pi)\le k<\rho\Delta\pi$: $t=0$
fails because $K_L(0)>k$; $t=\lambda$ fails because $K_L(\lambda)\le k$; an interior $t$ requires IM.
\end{proof}

\begin{lemma}[product structure; analytical]\label{lem:product}
Fix $r$ and $\pi$ with $\pi<\tau(r)$, and rates with $s\ge t$. Then $\mu_0(s,t)\le\pi<\tau$. The tuple with
rates $(s,t)$ is an equilibrium if and only if these four conditions hold:
\begin{enumerate}[label=(\roman*)]
\item $(s,e(1))$ is a pair of Lemma~\ref{lem:H};
\item $t$ is a rate of Lemma~\ref{lem:L};
\item $e(0)=e(-1)=\rho$;
\item the prices are $P(x)=t_0+e(x)[t_L-t_0+\Delta\mu(x)]$.
\end{enumerate}
\end{lemma}

\begin{proof}
If $s\ge t$, then $b_0+\lambda-s\le b_0+\lambda-t$, so the odds of $\mu_0$ are at most $\pi/(1-\pi)$. Hence
$\mu_0\le\pi<\tau$.

\emph{Only if.} Let the tuple be an equilibrium. Lemma~\ref{lem:price} applies to it. Part (d) gives
$e(0)=e(-1)=\rho$, since $\mu_0\le\pi<\tau$ and $\mu_-(t)\le\pi<\tau$; this is (iii). Part (b) gives (iv). Part (c)
makes the belief at the price of each order equal to the posterior of that order, so (E3) after the order $1$ holds
at the belief $\mu_+(s)$. By Lemma~\ref{lem:trade}(d), (E5) for class $H$ involves only $(s,\mu(1),e(1))$, and (E5)
for class $L$ only $(t,\mu(-1),e(-1))$. So (i) and (ii) hold.

\emph{If.} Suppose (i)--(iv), and set $\zeta(P(x))=(e(x)-\rho)/(1-\rho)$. The verification uses
Lemma~\ref{lem:price}(f), not parts (a)--(e).
\begin{itemize}
\item \emph{Entry is a nondecreasing function of the posterior.} The orders $0$ and $-1$ have posteriors at most
  $\pi<\tau$ and entry $\rho$. The order $1$ has the posterior $\mu_+(s)\ge\pi$ and entry $e(1)\ge\rho$. If
  $\mu_+(s)$ equals the posterior of another order, then $\mu_+(s)\le\pi$, so $s=0$; then (i) is pair (N) and
  $e(1)=\rho$. So orders with equal posteriors have equal entry, and a higher posterior never has lower entry.
\item \emph{Prices.} Lemma~\ref{lem:price}(f), with $\Delta>0$ and $t_L\ge t_0$, gives $P(x)=P(x')$ exactly when
  $\mu(x)=\mu(x')$. So $\zeta$ is well defined on the image of $P$.
\item (E1) holds by the flow law $f_H$, $f_L$.
\item (E2): the orders with the price $P(x)$ share the posterior $\mu(x)$, so the challenger's belief at $P(x)$ is
  their common posterior $\mu(x)$.
\item (E3): at the price of the order $1$, the belief is $\mu_+(s)$, and (i) states that $e(1)$ is optimal for it.
  At the prices of the orders $0$ and $-1$, the belief is below $\tau$, and (iii) gives $\zeta=0$.
\item (E4): entry after the order $x$ is $e(x)$, independent of $(j,\theta,R)$. So
  $\E[V_T\mid X=x]=t_0+e(x)[\mu(x)t_H+(1-\mu(x))t_L-t_0]$, which is (iv).
\item (E5): the tuple satisfies the premises of Lemma~\ref{lem:trade}. By its parts (a), (b) and (d), (E5) for class
  $H$ is the comparison of $K_H=e(1)\Delta[1-\mu_+(s)]$ with $k$ at the rate $s$, which (i) gives. (E5) for class
  $L$ is the comparison of $K_L=\rho\Delta\mu_-(t)$ with $k$ at the rate $t$, which (ii) gives. Entry after the zero
  order enters no payoff.
\item (E6) holds.\qedhere
\end{itemize}
\end{proof}

\subsection{The equilibrium set at prior $1/2$}

At $\pi=1/2$, $1-\mu_+(s)=b/(2b+s)$, $\mu_-(t)=b/(2b+t)$ and $s^*=b(2\tau-1)/(1-\tau)$. Five forms of the high
class's pair appear. In each, the low class uses the rate $\bar t(r)$ of Lemma~\ref{lem:L}:
\[
\bar t(r)=0\ \text{on}\ [r_0,r_N],\qquad \bar t(r)=b\,[\rho\Delta(r)/k-2]\ \text{on}\ (r_N,r_S),\qquad \bar t(r)=\lambda\ \text{on}\ [r_S,r_2].
\]
\begin{itemize}
\item \textbf{N} (no informed trade): $s=0$, $e(1)=\rho$.
\item \textbf{R} (a mixed purchase that recruits): $s=s_R(r):=b[\Delta(r)/k-2]$, $e(1)=1$.
\item \textbf{F} (a full purchase): $s=\lambda$; $e(1)=1$ for $r<r_C$, and $e(1)=\rho$ for $r>r_C$. At $r_C$,
  $e(1)=1$ under C1 and any $e(1)\in[\rho,1]$ under CM. The low class uses $\bar t(r)$, so on $(r_N,r_S)$ form F
  has a mixed sale: ``full purchase'' refers to the high class only.
\item \textbf{Q} (mixed orders at the entry floor): $s=s_Q(r):=b[\rho\Delta(r)/k-2]=\bar t(r)$, $e(1)=\rho$.
\item \textbf{T} (the threshold pair): $s=s^*(r)$, $e(1)=e^*(r)=k/J(r)$.
\end{itemize}
Each form defines the equilibrium tuple completely:
\begin{itemize}
\item orders: $\sigma_H(1)=s/\lambda$, $\sigma_H(0)=1-s/\lambda$, $\sigma_L(-1)=\bar t(r)/\lambda$, $\sigma_L(0)=1-\bar t(r)/\lambda$;
\item posteriors: $\mu_-(\bar t)$, $\mu_0(s,\bar t)$ and $\mu_+(s)$ after the orders $-1$, $0$ and $1$;
\item high-cost preparation: $(e(1)-\rho)/(1-\rho)$ at the price of the order $1$, and zero at the other prices;
\item prices: $P(x)=t_0+e(x)[t_L-t_0+\Delta\mu(x)]$.
\end{itemize}

\setcounter{proposition}{2}
\begin{proposition}[the equilibrium set along the benchmark path; computer-assisted]\label{prop:3}
Take the benchmark auction, the prior $1/2$ and any slot in the band. Under the rules IM and C1, the equilibrium
set at each strength is
\begin{center}
\begin{tabular}{@{}ll@{}}
\toprule
Strength & Equilibria\\
\midrule
$r_0\le r<r_A$ & N\\
$r=r_A$ & N; T with $e^*=1$ (the left end of R)\\
$r_A<r<r_F$ & N; R\\
$r_F\le r\le r_N$ & N; F\\
$r_N<r<r_J$ & F; Q\\
$r_J\le r\le r_2$ & F\\
\bottomrule
\end{tabular}
\end{center}
Under CM, T is also an equilibrium on $(r_A,r_J]$, and at $r=r_C$ entry after a purchase in F can be any
$e(1)\in[\rho,1]$. Under I0 the set is N on $[r_0,r_F]$; N and F on $(r_F,r_N]$; empty on $(r_N,r_S]$; and F on
$(r_S,r_2]$. Entry in each form is
\[
\Ent_N=\Ent_Q=\rho,\qquad \Ent_R(r)=\rho+(1-\rho)\frac{b\,\Delta(r)}{2k},
\]
\[
\Ent_F=\begin{cases}\Ent_1=\rho+(1-\rho)(b+\lambda/2), & r<r_C;\ \text{also } r=r_C \text{ under C1},\\
\rho+[e(1)-\rho](b+\lambda/2)\in[\rho,\Ent_1], & r=r_C \text{ under CM},\ e(1)\in[\rho,1],\\
\rho, & r>r_C,\end{cases}
\]
\[
\Ent_T(r)=\rho+\frac b2\cdot\frac{k-\rho J(r)}{\Delta(r)[1-\tau(r)]^2}.
\]
Hence:
\begin{enumerate}[label=(\alph*)]
\item No informed trade and an informative equilibrium coexist exactly on $[r_A,r_N]$ (on $(r_F,r_N]$ under I0).
\item On the informative branch that coexists with no trade, entry is $\rho+(1-\rho)\min\{b\Delta(r)/(2k),\,b+\lambda/2\}$.
  It is continuous, rises strictly on $[r_A,r_F]$, and is constant at $\Ent_1$ on $[r_F,r_N]$. It stays at
  $\Ent_1$ along F up to $r_C$ and equals $\rho$ above $r_C$. At $r_C$ itself, entry is $\Ent_1$ under C1; under CM
  it can take every value in $[\rho,\Ent_1]$. The strict rise needs IM.
\item Under CM, entry along T falls strictly from $\Ent_R(r_A)$ at $r_A$ to $\rho$ at $r_J$. For
  $r_A<r\le r_N$, $\rho<\Ent_T(r)<\min\{\Ent_R(r),\Ent_1\}$.
\item The equilibrium set is nonempty at every strength under IM, with either challenger rule.
\end{enumerate}
\end{proposition}

\begin{proof}
\emph{Step 1: the frame.} By Lemma~\ref{lem:path}(a), $\tau(r)\ge\tau(r_0)>1/2=\pi$ on the path, so
Lemmas~\ref{lem:H}--\ref{lem:product} apply at every $r$. Every equilibrium has a high-class pair from
Lemma~\ref{lem:H} and a low-class rate from Lemma~\ref{lem:L}. Step 4 shows $s\ge t$ for every such combination,
so Lemma~\ref{lem:product} gives the equilibrium set as their product.

\emph{Step 2: the high-class pairs.} Use Lemma~\ref{lem:path}, Definition~\ref{def:thresholds} and
Lemma~\ref{lem:order}. Recall $1-M=1/(2+u)$ and $ub=\lambda$.
\begin{itemize}
\item (N) holds iff $\rho\Delta/2\le k$, that is $r\le r_N$.
\item (R): $\Delta b/(2b+s)=k$ gives $s=s_R(r)$. Then $\mu_+(s_R)>\tau$ iff $k/\Delta<1-\tau$ iff $J(r)>k$ iff
  $r>r_A$. And $s_R<\lambda$ iff $\Delta/k<2+u$ iff $\Delta(1-M)<k$ iff $r<r_F$. So (R) holds iff $r_A<r<r_F$,
  under IM.
\item (Q): $\rho\Delta b/(2b+s)=k$ gives $s=s_Q(r)$. Then $s_Q>0$ iff $r>r_N$; $\mu_+(s_Q)<\tau$ iff $\rho J(r)<k$
  iff $r<r_J$; $s_Q<\lambda$ iff $\rho\Delta(1-M)<k$ iff $r<r_S$. Since $r_J<r_S$, (Q) holds iff $r_N<r<r_J$,
  under IM.
\item (T): $s^*\in(0,\lambda)$ iff $1/2<\tau(r)<M$ iff $r<r_C$. Also $\rho\le e^*\le1$ iff $\rho J(r)\le k\le J(r)$
  iff $r_A\le r\le r_J$, and $e^*=1$ iff $r=r_A$. So (T) holds at $r_A$ under C1, and on $[r_A,r_J]$ under CM,
  under IM in both cases.
\item (F): for $r<r_C$, $M>\tau(r)$, so $e(1)=1$, and the condition is $\Delta(1-M)\ge k$, that is $r\ge r_F$
  ($r>r_F$ under I0). For $r>r_C$, $M<\tau(r)$, so $e(1)=\rho$, and the condition $\rho\Delta(1-M)\ge k$ holds
  because $r>r_C>r_S$. At $r=r_C$, C1 gives $e(1)=1$; CM allows any $e(1)\in[\rho,1]$, and every such value
  satisfies the condition because $\rho\Delta(r_C)(1-M)>k$. So (F) holds iff $r\ge r_F$ (IM) or $r>r_F$ (I0).
\end{itemize}

\emph{Step 3: the low-class rate.} At $\pi=1/2$, Lemma~\ref{lem:L} gives $t=0$ iff $r\le r_N$, $t=\lambda$ iff
$r\ge r_S$, and $\bar t(r)=b(\rho\Delta/k-2)$ on $(r_N,r_S)$, under IM. Under I0 no rate qualifies iff
$\rho\Delta m\le k<\rho\Delta/2$, that is on $(r_N,r_S]$.

\emph{Step 4: $s\ge t$.} For $r\le r_N$, $t=0\le s$. For $r>r_N$, the pairs are F with $s=\lambda\ge t$; Q with
$s=s_Q(r)=\bar t(r)$; and T with $r_N<r\le r_J$. For T, $s^*\ge \bar t(r)$ iff $(2\tau-1)/(1-\tau)\ge\rho\Delta/k-2$ iff
$1/(1-\tau)\ge\rho\Delta/k$ iff $\rho J(r)\le k$, which holds for $r\le r_J$. By Lemma~\ref{lem:product}, entry
after the zero order and after a sale is $\rho$ in every equilibrium.

\emph{Step 5: the table.} Combine Steps 2--4 with the order $r_A<r_F<r_N<r_J<r_S<r_C$ of Lemma~\ref{lem:order}.
On $[r_J,r_2]$ only F remains. On $(r_N,r_J)$, F and Q remain. On $[r_F,r_N]$, N and F remain, and so on. Under I0,
(Q), (R) and (T) drop out, (F) needs $r>r_F$, and Step 3 removes every equilibrium on $(r_N,r_S]$.

\emph{Step 6: verification of each tuple.} Fix a form and a strength at which Steps 2--3 admit it, and take the
tuple that the form defines. Its rates satisfy $s\ge t$ by Step 4. Condition (i) of Lemma~\ref{lem:product} holds by
Step 2, condition (ii) by Step 3, and conditions (iii) and (iv) by the definition of the tuple. The ``if'' direction
of Lemma~\ref{lem:product} then verifies (E1)--(E6) directly. It derives the price separation from the price
formula through Lemma~\ref{lem:price}(f), and it assumes no equilibrium condition.

\emph{Step 7: entry.} Entry is $\Ent=\sum_x\Pr(X=x)e(x)$. Since $e(0)=e(-1)=\rho$ and
$\Pr(X=1)=\tfrac12(b+s)+\tfrac12b=b+s/2$, $\Ent=\rho+[e(1)-\rho](b+s/2)$.
\begin{itemize}
\item R: $b+s_R/2=b\Delta/(2k)$ and $e(1)=1$. This is strictly increasing in $r$ because $\Delta$ is. At $r_F$,
  $\Delta(r_F)=k/(1-M)=k(2+u)$, so $b\Delta(r_F)/(2k)=b+ub/2=b+\lambda/2$: the branch joins F continuously.
\item F: $b+\lambda/2$, with $e(1)=1$ below $r_C$ and $e(1)=\rho$ above. At $r_C$, $e(1)=1$ under C1, and
  under CM every $e(1)\in[\rho,1]$ gives an equilibrium (Step 2), so entry fills $[\rho,\Ent_1]$.
\item Q: $e(1)=\rho$, so $\Ent=\rho$.
\item T: $e^*-\rho=(k-\rho J)/J$ and $b+s^*/2=b/[2(1-\tau)]$, so $\Ent_T-\rho=\tfrac b2(k-\rho J)/[J(1-\tau)]$,
  which is the stated formula because $J(1-\tau)=\Delta(1-\tau)^2$.
\end{itemize}
At $r_A$, $J=k$ and $\Delta(r_A)/k=1/(1-\tau(r_A))$, so $\Ent_T(r_A)=\rho+\tfrac12(1-\rho)b/(1-\tau(r_A))=\Ent_R(r_A)$.
At $r_J$, $\rho J=k$, so $\Ent_T(r_J)=\rho$. On $[r_A,r_J]$, $k-\rho J$ is positive except at $r_J$ and strictly
decreasing (Lemma~\ref{lem:path}(b)). Next, $\rho J(2)-k=11/11200>0$, so $r_J<2$, and
Lemma~\ref{lem:path}(c) makes $\Delta(1-\tau)^2$ positive and strictly increasing on $[r_A,r_J]\subset[r_0,2]$.
So $\Ent_T$ falls strictly. For $r_A<r\le r_N$: $r<r_J$ gives $e^*>\rho$, so
$\Ent_T>\rho$; $J>k$ gives $e^*<1$, and $s^*<s$ in R or F, so $\Ent_T<\Ent_R$ or $\Ent_T<\Ent_1$.

\emph{Step 8: existence.} Under IM, N exists on $[r_0,r_N]$ and F on $[r_F,r_2]$, and $r_F<r_N$.
\end{proof}

\paragraph{Calibration B.} At calibration B the table reads as follows (\lab{computer-assisted}):
\begin{center}
\small
\begin{tabular}{@{}lll@{}}
\toprule
Strengths & Equilibria (IM, C1) & Entry\\
\midrule
$[1.2,\,1.393958)$ & N & $0.250000$\\
$1.393958$ & N; T $=$ left end of R & $0.250000$; $0.454185$\\
$(1.393958,\,1.467127)$ & N; R & $0.250000$; from $0.454185$ up to $0.522757$\\
$[1.467127,\,1.747878]$ & N; F & $0.250000$; $0.522757$\\
$(1.747878,\,1.966097)$ & F; Q & $0.522757$; $0.250000$\\
$[1.966097,\,3.592659]$ & F & $0.522757$\\
$(3.592659,\,3.6]$ & F & $0.250000$\\
\bottomrule
\end{tabular}
\end{center}
At $r_C=3.592659$ the table gives entry $0.522757$ under C1. Under CM, entry at $r_C$ can be any value in
$[0.250000,\,0.522757]$.
On branch R, entry is $0.459588$ at $1.40$ and $0.506112$ at $1.45$. The low class abstains on
$[1.2,1.747878]$, mixes on $(1.747878,2.124148)$ and sells fully on $[2.124148,3.6]$. The strengths $1.55$, $1.60$
and $1.65$ lie in $[r_F,r_N]$, where the equilibria are N and F under C1. Under CM the threshold pair T adds
entry $0.341462$, $0.321311$ and $0.305208$ at these three strengths. The script \texttt{p1\_enumerate.py} rebuilds
the equilibrium set from the flow law at 241 strengths and all four pairs of rules, at calibrations A and B.
It finds no difference from the table (\lab{numerical diagnostic} at A in exact arithmetic, at B in interval
arithmetic).

\subsection{Coexistence at every prior}

For a prior $\pi$, write $r_N(\pi)=\fr\big(k/(\rho\max\{\pi,1-\pi\})\big)$. The two other cut points exist only on
part of the prior range, and the proof uses each one only there.
\begin{itemize}
\item \emph{The ceiling $r_C(\pi)$.} By Lemma~\ref{lem:path}(a), $\{r\in[r_0,r_2]:\tau(r)\le M(\pi)\}$ is empty when
  $M(\pi)<\tau(r_0)$, and otherwise it is an interval $[r_0,r_C(\pi)]$. Define $r_C(\pi)$ only when
  $M(\pi)\ge\tau(r_0)$.
\item \emph{The prior tie $r_\pi$.} For $\tau(r_0)\le\pi\le\tau(r_2)$, let $r_\pi\in[r_0,r_2]$ solve
  $\tau(r_\pi)=\pi$; it is unique because $\tau$ is continuous and strictly increasing. For $\pi>\tau(r_2)$, no
  strength on the path has $\pi<\tau(r)$, and $r_\pi$ is not defined.
\end{itemize}
Write $\tau_A=\tau(r_A)$, and let $r_*$ solve $\rho\Delta(r_*)\tau(r_*)=k$ on the path (unique by
Lemma~\ref{lem:path}(d)). Define
\[
\pi_A^{\rm lo}=\frac{\tau_A}{\tau_A+(1-\tau_A)(1+u)},\qquad \pi_A^{\rm hi}=\tau_A,\qquad \pi_{\max}=\tau(r_*).
\]
So $M(\pi_A^{\rm lo})=\tau_A$. The numbers $\pi_A^{\rm hi}$ and $\pi_{\max}$ do not depend on the slot. Since
$r_0<r_A<r_*<r_2$, they satisfy $\tau(r_0)<\pi_A^{\rm hi}<\pi_{\max}<\tau(r_2)$; at the benchmark
$\pi_{\max}-\tau_A=0.009320$. In cases (b) and (c) below, $M(\pi)>\tau_A>\tau(r_0)$, so $r_C(\pi)$ exists; in
case (c), $\pi\in[\tau_A,\pi_{\max})\subset[\tau(r_0),\tau(r_2)]$, so $r_\pi$ exists.

\setcounter{proposition}{4}
\begin{proposition}[coexistence at every prior; computer-assisted]\label{prop:5}
Take the benchmark auction, any slot with $0<\lambda<1$ and $0<\beta<1/2$, any prior $\pi\in(0,1)$ and the rules IM
and C1. The set $C(\pi)$ of strengths on the path at which no informed trade and an informative equilibrium
coexist is
\[
C(\pi)=\big\{r\in[r_0,r_2]:\ r\ge r_A,\ \ \pi<\tau(r)\le M(\pi),\ \ \rho\Delta(r)\max\{\pi,1-\pi\}\le k\big\}.
\]
\begin{enumerate}[label=(\alph*)]
\item If $\pi<\pi_A^{\rm lo}$, $C(\pi)$ is empty. At $\pi=\pi_A^{\rm lo}$, $C(\pi)=\{r_A\}$.
\item If $\pi_A^{\rm lo}<\pi<\pi_A^{\rm hi}$, then $C(\pi)=[r_A,\min\{r_N(\pi),r_C(\pi)\}]$, a nondegenerate interval.
  At $r_A$ the informative equilibrium is the threshold pair: $s=s^*$, $e(1)=1$, and both the investor and the
  challenger are indifferent. The start $r_A$ is the same for every prior in this window and every slot.
\item If $\pi_A^{\rm hi}\le\pi<\pi_{\max}$, then $C(\pi)=(r_\pi,\min\{r_N(\pi),r_C(\pi)\}]$, a nonempty interval
  with $r_\pi\ge r_A$. Coexistence starts at $r_\pi$, the strength at which the prior reaches the threshold.
\item If $\pi\ge\pi_{\max}$, $C(\pi)$ is empty.
\end{enumerate}
At each $r\in C(\pi)$ the equilibrium set is N and one informative equilibrium: R with the rate $s$ that solves
$\Delta[1-\mu_+(s)]=k$ when $\Delta(1-M(\pi))<k$, and F when $\Delta(1-M(\pi))\ge k$ (T at $r_A$ in case (b)).
Its entry is $\rho+(1-\rho)\min\{(1-\pi)b\Delta(r)/k,\ b+\pi\lambda\}$. It rises strictly while the purchase is mixed
and then stays constant. At calibration B, $\pi_A^{\rm lo}=0.396175$, $\pi_A^{\rm hi}=0.640739$ and
$\pi_{\max}=0.650059$.
\end{proposition}

\begin{proof}
\emph{Step 1: strengths with $\pi<\tau(r)$.} Lemmas~\ref{lem:H}--\ref{lem:L} apply. No informed trade is an
equilibrium if and only if (N) holds in both lemmas, that is $\rho\Delta\max\{\pi,1-\pi\}\le k$; then
$\mu_0=\pi<\tau$ and $e(0)=\rho$. Suppose this holds. Lemma~\ref{lem:L} then gives $t=0$ in every equilibrium, so
an informative equilibrium has $s>0$. Pair (Q) is impossible: $\rho\Delta[1-\mu_+(s)]<\rho\Delta(1-\pi)\le k$ for
$s>0$. Pair (F) with $M(\pi)<\tau$ is impossible for the same reason. The other pairs, (R), (T) and (F) with
$M(\pi)\ge\tau$, need $\mu_+(s)\ge\tau$ for some $s\le\lambda$, so $\tau\le M(\pi)$. They also need
$k\le e(1)\Delta[1-\mu_+(s)]\le\Delta(1-\tau)=J(r)$, so $r\ge r_A$ by Lemma~\ref{lem:path}(b).

Conversely, let $\tau(r)\le M(\pi)$ and $r\ge r_A$, so $J(r)\ge k$. If $\Delta(1-M(\pi))\ge k$, pair (F) holds with
$e(1)=1$. Otherwise $\tau<M(\pi)$, and $s\mapsto\Delta[1-\mu_+(s)]$ is continuous and strictly decreasing on
$[s^*,\lambda]$, from $J(r)\ge k$ to $\Delta(1-M(\pi))<k$. It equals $k$ at exactly one $s\in[s^*,\lambda)$. At
$s=s^*$ this is (T) with $J(r)=k$ and $e(1)=1$; at $s>s^*$ it is (R). With $t=0\le s$, Lemma~\ref{lem:product}
completes the tuple. So a strength with $\pi<\tau(r)$ lies in $C(\pi)$ if and only if it satisfies the three
conditions in the display.

\emph{Step 2: strengths with $\pi\ge\tau(r)$.} Under C1, no informed trade gives the posterior $\pi\ge\tau$ after
every order, so $e\equiv1$, and it is an equilibrium if and only if $\Delta\max\{\pi,1-\pi\}\le k$. Suppose this
holds, and take any candidate with $s>0$. Then $\mu_+(s)>\pi\ge\tau$, so $e(1)=1$ and
$K_H=\Delta[1-\mu_+(s)]<\Delta(1-\pi)\le k$: class $H$ strictly prefers abstention, a contradiction. Take any
candidate with $t>0$. Then $\mu_-(t)<\pi$ and $e(-1)\le1$, so $K_L\le\Delta\mu_-(t)<\Delta\pi\le k$: again a
contradiction. So no strength with $\pi\ge\tau(r)$ lies in $C(\pi)$.

\emph{Step 3: the four cases.} By Lemma~\ref{lem:path}(a), the condition $\pi<\tau(r)$ holds
\begin{itemize}
\item for every strength on the path if $\pi<\tau(r_0)$;
\item exactly for $r>r_\pi$ if $\tau(r_0)\le\pi\le\tau(r_2)$ (at $\pi=\tau(r_0)$, $r_\pi=r_0$ and the strict
  inequality holds only above $r_0$);
\item for no strength if $\pi>\tau(r_2)$.
\end{itemize}
The condition $\tau(r)\le M(\pi)$ holds exactly for $r\le r_C(\pi)$ when $M(\pi)\ge\tau(r_0)$, and for no strength
otherwise. The last condition holds for $r\le r_N(\pi)$.
Also $\rho\Delta(r_A)\max\{\pi,1-\pi\}<\rho\Delta(r_A)=\rho k/(1-\tau_A)<k$, because $\tau_A<3/4=1-\rho$ (the margin
$3/4-\tau_A=0.109261$ is a derived strict inequality of Table~\ref{tab:derived}); so $r_A<r_N(\pi)$ for every prior.
\begin{itemize}
\item[(a)] $M$ is increasing in $\pi$, so $\pi<\pi_A^{\rm lo}$ gives $M(\pi)<\tau_A\le\tau(r)$ for every $r\ge r_A$.
  No strength qualifies. At $\pi=\pi_A^{\rm lo}$ only $r_A$ qualifies.
\item[(b)] For $r\ge r_A$, $\tau(r)\ge\tau_A>\pi$. Also $\tau_A<M(\pi)$ gives $r_A<r_C(\pi)$, and $r_A<r_N(\pi)$.
  At $r_A$, $\Delta(r_A)(1-M(\pi))<\Delta(r_A)(1-\tau_A)=k$, so Step 1 gives the threshold pair with $J=k$.
  There $\mu_+=\tau$, so the challenger is indifferent, and $K_H=k$, so the investor is indifferent.
\item[(c)] Now $r_\pi\ge r_A$, and $\pi<\tau(r)$ holds exactly for $r>r_\pi$. The interval is nonempty: $r_\pi<r_C(\pi)$
  because $\tau(r_\pi)=\pi<M(\pi)$; and $r_\pi<r_N(\pi)$ iff $\rho\Delta(r_\pi)\pi<k$ (as $\pi>1/2$) iff
  $\rho\Delta(r_\pi)\tau(r_\pi)<k$ iff $r_\pi<r_*$ (Lemma~\ref{lem:path}(d)) iff $\pi<\pi_{\max}$.
\item[(d)] Here $\pi>1/2$. If $\pi>\tau(r_2)$, no strength has $\pi<\tau(r)$, and Step 2 covers every strength.
  Otherwise $\pi_{\max}\le\pi\le\tau(r_2)$, so $r_\pi$ exists and $r_\pi\ge r_*$. Any $r$ with $\pi<\tau(r)$ has
  $r>r_\pi\ge r_*$, so $\rho\Delta(r)\pi>\rho\Delta(r_*)\tau(r_*)=k$, and N is not an equilibrium there. Step 2
  covers the other strengths.
\end{itemize}
The entry formula follows as in Step 7 of Proposition~\ref{prop:3}: $\Pr(X=1)=b+\pi s$, and in R the equation
$\Delta(1-\pi)b/(b+\pi s)=k$ gives $b+\pi s=(1-\pi)b\Delta/k$; in F, $\Pr(X=1)=b+\pi\lambda$.
\end{proof}

\paragraph{Other tie rules.} Write $C(\pi)$ for the set of Proposition~\ref{prop:5} (rules IM and C1). The other
three rule pairs change it only at a prior tie $\tau(r)=\pi$ and through the strict full purchase of I0.

\begin{remark}[coexistence under the other tie rules; analytical]\label{rem:ties}
Take the setting of Proposition~\ref{prop:5}. Let $C_{\rm CM}(\pi)$, $C_{\rm I0,C1}(\pi)$ and $C_{\rm I0,CM}(\pi)$ be
the coexistence sets under the rule pairs (IM, CM), (I0, C1) and (I0, CM).
\begin{enumerate}[label=(\alph*)]
\item $C_{\rm CM}(\pi)=C(\pi)\cup\{r_\pi\}$ if $\tau_A<\pi\le\pi_{\max}$, and $C_{\rm CM}(\pi)=C(\pi)$ otherwise. The
  lower end $\pi=\tau_A$ is excluded, and the upper end $\pi=\pi_{\max}$ is included.
\item $C_{\rm I0,C1}(\pi)=\{r\in C(\pi):\Delta(r)[1-M(\pi)]>k\}$, the strengths of $C(\pi)$ above
  $r_F(\pi)=\fr\big(k/(1-M(\pi))\big)$.
\item $C_{\rm I0,CM}(\pi)=C_{\rm I0,C1}(\pi)\cup\{r_\pi\}$ if $\tau(r_0)\le\pi\le\pi_{\max}$ and
  $\Delta(r_\pi)[1-M(\pi)]>k$, and $C_{\rm I0,CM}(\pi)=C_{\rm I0,C1}(\pi)$ otherwise.
\end{enumerate}
In particular, at $\pi=\pi_{\max}$ the set $C(\pi)$ is empty, but $C_{\rm CM}(\pi)=\{r_*\}$. At calibration B,
$\Delta(r_*)[1-M(\pi_{\max})]-k=0.000343>0$, so also $C_{\rm I0,CM}(\pi_{\max})=\{r_*\}$, while
$C_{\rm I0,C1}(\pi_{\max})$ is empty.
\end{remark}

\begin{proof}
\emph{Step 1: strengths with $\tau(r)\ne\pi$.} If $\pi>\tau(r)$, every posterior of the no-trade tuple is
$\pi>\tau$, so entry is $1$ under both challenger rules. Step 2 of Proposition~\ref{prop:5} uses only strict
comparisons, so it excludes coexistence under every rule pair. If $\pi<\tau(r)$, no-trade posteriors are
$\pi<\tau$, with no tie. Step 1 of Proposition~\ref{prop:5} shows that an informative pair compatible with no trade
needs $\tau\le M(\pi)$ and $J(r)\ge k$. Under CM, pair (T) needs $e^*=k/J\le1$ and $s^*\le\lambda$, and pair (F) at
$\tau=M(\pi)$ can use $e(1)=1$; so CM adds no strength. Under I0, pairs (Q), (R) and (T) need a mixture at
indifference and drop out. Pair (F) needs $e(1)\Delta[1-M(\pi)]>k$ with $e(1)\le1$, so $\Delta[1-M(\pi)]>k$; then
$\tau\le M(\pi)$ gives $e(1)=1$ under C1, and CM allows $e(1)=1$. No trade needs only weak inequalities, which I0
satisfies by abstention. This gives (b) and the part of (a) and (c) away from $r_\pi$.

\emph{Step 2: the prior tie $r=r_\pi$ under C1.} Here $\pi=\tau(r_\pi)\ge\tau(r_0)>1/2$. Step 2 of
Proposition~\ref{prop:5} covers $\pi\ge\tau(r)$, so $r_\pi\notin C(\pi)$ and $r_\pi\notin C_{\rm I0,C1}(\pi)$.

\emph{Step 3: the prior tie $r=r_\pi$ under CM.} Write $\Delta=\Delta(r_\pi)$.
\begin{itemize}
\item \emph{No trade.} Every order has the posterior $\pi=\tau$. CM allows $\zeta=0$ at the common price, so entry
  $\rho$ after every order is admissible, and it is the entry that best supports abstention. By
  Lemma~\ref{lem:trade}(b), no trade is an equilibrium if and only if $\rho\Delta\max\{\pi,1-\pi\}=\rho\Delta\pi\le k$.
  By Lemma~\ref{lem:path}(d), $\rho\Delta(r)\tau(r)$ increases strictly, so this holds if and only if $r_\pi\le r_*$,
  that is $\pi\le\pi_{\max}$.
\item \emph{The low class.} Suppose $\rho\Delta\pi\le k$. A rate $t>0$ gives $\mu_-(t)<\pi=\tau$, so $e(-1)=\rho$ and
  $K_L=\rho\Delta\mu_-(t)<\rho\Delta\pi\le k$. So $t=0$. After a sale the posterior is then $\pi=\tau$, and CM allows
  $e(-1)=\rho$, with $K_L=\rho\Delta\pi\le k$.
\item \emph{The high class.} A rate $s>0$ gives $\mu_+(s)>\tau$, so $e(1)=1$ and
  $K_H(s)=\Delta[1-\mu_+(s)]<\Delta(1-\pi)=J(r_\pi)$. An informative equilibrium therefore needs $J(r_\pi)>k$, that is
  $r_\pi>r_A$, that is $\pi>\tau_A$. Conversely, let $J(r_\pi)>k$. The return $K_H$ falls continuously on $(0,\lambda]$
  from the limit $J(r_\pi)>k$ to $\Delta[1-M(\pi)]$. If $\Delta[1-M(\pi)]\ge k$, the full purchase $s=\lambda$
  qualifies (strictly when $>k$). Otherwise one interior rate solves $K_H=k$ (rule IM only).
\item \emph{The tuple.} With $s>0=t$, the zero-order posterior is $\mu_0<\pi<\mu_+(s)$, and $\mu(-1)=\pi$. Take
  $e=(\rho,\rho,1)$, $P(x)=t_0+e(x)[t_L-t_0+\Delta\mu(x)]$, $\zeta=0$ at $P(-1)$ and $P(0)$, $\zeta=1$ at $P(1)$, and
  truthful bids. Entry is nondecreasing in the posterior, and the three posteriors differ, so
  Lemma~\ref{lem:price}(f) gives three distinct prices. (E2) gives the belief $\mu(x)$ at $P(x)$. (E3) holds: at
  $P(-1)$ the belief is $\tau$ and CM allows $\zeta=0$; at $P(0)$ it is below $\tau$; at $P(1)$ it is above $\tau$.
  (E4) holds by direct computation, as in the proof of Lemma~\ref{lem:product}. (E5) holds by the two class bullets,
  through Lemma~\ref{lem:trade}. (E6) holds.
\end{itemize}
So under (IM, CM), $r_\pi$ joins the coexistence set exactly when $\tau_A<\pi\le\pi_{\max}$. Under (I0, CM) the
informative tuple must be the strict full purchase, so $r_\pi$ joins exactly when $\pi\le\pi_{\max}$ and
$\Delta(r_\pi)[1-M(\pi)]>k$; this strict inequality implies $J(r_\pi)>k$. At $\pi=\pi_{\max}$, $r_\pi=r_*$, and the
low class is indifferent between a sale and abstention in both tuples; IM and I0 both allow abstention. The value
$0.000343$ at calibration B is the registry margin \texttt{margin\_cmend\_B}.
\end{proof}

Table~\ref{tab:cmend} lists the two tuples at the endpoint $r=r_*$, $\pi=\pi_{\max}$ at calibration B
(\lab{computer-assisted}: the values enclose the root bracket of $r_*$ in interval arithmetic).

\begin{table}[h]
\centering
\caption{The two equilibria at the CM endpoint $r=r_*=1.634218$, $\pi=\pi_{\max}=0.650059$, calibration B.}
\label{tab:cmend}
\small
\begin{tabular}{@{}lll@{}}
\toprule
Component & No trade & Informative trade\\
\midrule
Orders $(\sigma_H,\sigma_L)$ & $(\delta_0,\delta_0)$ & $(\delta_1,\delta_0)$\\
Posteriors after $(-1,0,1)$ & $(\pi,\pi,\pi)$ & $(\pi,\ 0.454140,\ 0.834698)$\\
Entry after $(-1,0,1)$ & $(\rho,\rho,\rho)$ & $(\rho,\rho,1)$\\
Prices & $t_0+\rho[t_L-t_0+\Delta\pi]$ at every order & $t_0+e(x)[t_L-t_0+\Delta\mu(x)]$\\
Class $H$: net return of a purchase & $\rho\Delta(1-\pi)-k<0$ & $\Delta[1-M(\pi)]-k=0.000343>0$\\
Class $L$: net return of a sale & $\rho\Delta\pi-k=0$ (tie) & $\rho\Delta\pi-k=0$ (tie)\\
Entry & $0.250000$ & $0.560586$\\
\bottomrule
\end{tabular}
\end{table}

The script \texttt{p1\_enumerate.py} also confirms part (a) away from the endpoint, in exact arithmetic: at the
grid strength $3/2$ and the prior $\tau(3/2)=69/107\in(\tau_A,\pi_{\max})$, coexistence holds under CM and fails
under C1 (\lab{numerical diagnostic}).

\subsection{Which conclusions depend on the tie rules}

\begin{center}
\small
\begin{tabular}{@{}lllll@{}}
\toprule
Result & IM, C1 & IM, CM & I0, C1 & I0, CM\\
\midrule
Proposition 2 (i)--(iii) & holds & holds & holds & holds\\
No informed trade & $[r_0,r_N]$ & $[r_0,r_N]$ & $[r_0,r_N]$ & $[r_0,r_N]$\\
Mixed purchase R & $(r_A,r_F)$ & $(r_A,r_F)$ & none & none\\
Threshold pair T & $r_A$ only & $[r_A,r_J]$ & none & none\\
Full purchase F & $[r_F,r_2]$ & $[r_F,r_2]$ & $(r_F,r_N]\cup(r_S,r_2]$ & $(r_F,r_N]\cup(r_S,r_2]$\\
Equilibrium set on $(r_N,r_J)$ & F; Q & F; Q; T & empty & empty\\
Equilibrium set at $r_J$ & F & F; T & empty & empty\\
Equilibrium set on $(r_J,r_S]$ & F & F & empty & empty\\
Entry after a purchase at $r_C$ & $1$ & any in $[\rho,1]$ & $1$ & any in $[\rho,1]$\\
Coexistence at prior $1/2$ & $[r_A,r_N]$ & $[r_A,r_N]$ & $(r_F,r_N]$ & $(r_F,r_N]$\\
Coexistence at prior $\pi$ & $C(\pi)$ & Remark~\ref{rem:ties}(a) & Remark~\ref{rem:ties}(b) & Remark~\ref{rem:ties}(c)\\
\bottomrule
\end{tabular}
\end{center}
Under I0, F on $(r_N,r_S]$ is not an equilibrium: the low class has no admissible rate there. The investor rule
IM is therefore the rule under which the equilibrium set is nonempty along the whole path. The rise of entry along
R also needs IM. The challenger rule matters only through T, the point $r_C$, and (at priors above
$\tau(r_0)$) the prior tie $r_\pi$.

\section{The admissible region and the margins table}\label{sec:region}

\setcounter{proposition}{3}
\begin{proposition}[the admissible region; analytical]\label{prop:region}
Take the benchmark primitives and strengths, and a slot with $0<\lambda<1$ and $0<\beta\le1/3$.
\begin{enumerate}[label=(\alph*)]
\item At the prior $1/2$, (H1)--(H3) hold if and only if
  \[
  \frac{82}{59}<u<\frac{2686}{1559},\qquad\text{that is}\qquad \tau(r_1)=\frac{141}{200}<M<\frac{4245}{5804}=\tau(r_2).
  \]
  In terms of the slot, $82\beta/(59+82\beta)<\lambda<2686\beta/(1559+2686\beta)$.
\item On each side of the band a conclusion of Proposition~\ref{prop:2} fails. If $u<82/59$, entry at $r_1$ is
  $\rho$. If $u=82/59$, entry at $r_1$ is $\Ent_1$ under C1 and can be $\rho$ under CM. If $2686/1559\le u<19/3$,
  entry after a purchase at $r_2$ is $1$ under C1. For $u\ge19/3$, the right side of (H2) fails at $r_1$.
\item At calibration B, with the other benchmark inputs fixed, (H1)--(H3) hold if and only if
  \[
  \pi_{\rm lo}:=\frac{141}{141+59e}<\pi<\frac{4245}{4245+1559e}=:\pi_{\rm hi},\qquad
  \pi_{\rm lo}\approx0.467850,\quad \pi_{\rm hi}\approx0.500424.
  \]
  Hypotheses (H1)--(H2), which give parts (i)--(ii), hold if and only if $\pi_{\rm lo}<\pi<195/308$, with
  $195/308\approx0.633117$.
\end{enumerate}
\end{proposition}

\begin{proof}
(a) At $\pi=1/2$, $\kappa=m=1-M=1/(2+u)$. The left side of (H1) is $c_H-B_{r_0}(1/2)=287/240>0$, and the left
side of (H2) is $k-\Delta(r_0)=1/300>0$; neither depends on the slot. The right side of (H1) is
$M>\tau(r_1)=141/200$, that is $200(1+u)>141(2+u)$, that is $u>82/59$. Hypothesis (H3) is $M<4245/5804$, that is
$5804(1+u)<4245(2+u)$, that is $u<2686/1559$. The right side of (H2) is $(1-M)\Delta(r_1)/4>1/50$, that is
$M<22/25$, that is $u<19/3$; the bound $u<2686/1559$ implies it. Finally $\lambda=u\beta/(1+u\beta)$ is increasing
in $u$ at fixed $\beta$.

(b) For $u<19/3$ the bound of Step 2 of Proposition~\ref{prop:2} holds at $r_1$ and at $r_2$, since
$\rho(1-M)\Delta(r_2)>k$ holds for $M<773/845$. So orders are full at both strengths and the posteriors are
$(m,1/2,M)$. If $M<\tau(r_1)$, no order recruits at $r_1$. If $M=\tau(r_1)$, a purchase leaves the challenger
indifferent, and the tie rule decides. If $M\ge\tau(r_2)$, a purchase recruits at $r_2$ under C1.

(c) $M(\pi)>\tau$ is $\pi/(1-\pi)>\tau/[(1-\tau)(1+u)]$. With $1+u=e$, the right side of (H1) gives
$\pi>141/(141+59e)$, and (H3) gives $\pi<4245/(4245+1559e)$. The left side of (H1) is $\pi<\tau(r_0)=195/308$. The
right side of (H2) is $\kappa(\pi)>3/25$: $m(\pi)>3/25$ is $\pi>3e/(22+3e)\approx0.270432$, and $1-M(\pi)>3/25$ is
$\pi<22/(22+3e)\approx0.729568$. The left side of (H2) does not depend on $\pi$. Intersect the intervals.
\end{proof}

The same band carries Proposition~\ref{prop:3}: the proof of Lemma~\ref{lem:order} uses only
$\tau(r_1)<M<\tau(r_2)$, and Proposition~\ref{prop:3} uses only Lemmas~\ref{lem:flow}(a),(b),(d)
and~\ref{lem:order}.

\paragraph{The open parameter set.} The admissible slots form the set
$\{(\lambda,\beta):0<\beta<1/2,\ 82/59<\lambda/[(1-\lambda)\beta]<2686/1559\}$, which is open. The model restriction
$\beta\le1/3$ is a structural weak restriction, not a strict hypothesis; it selects a relatively open section of
this set. At $\beta=1/3$ the band is $82/259<\lambda<2686/7363$, that is approximately $(0.316602,0.364797)$. At
$\beta=0.294655$ it is $\lambda_{\rm lo}(\beta)<\lambda<\lambda_{\rm hi}(\beta)$ with
$\lambda_{\rm lo}(\beta)=82\beta/(59+82\beta)=2416171/8316171\approx0.290539$ and
$\lambda_{\rm hi}(\beta)=2686\beta/(1559+2686\beta)=79144333/235044333\approx0.336721$; calibration B has
$\lambda\approx0.336122$. For parts (i)--(ii) alone the slot's prior range is wider than the Laplace range
$(0.467850,0.538096)$, because the slot's trading bound $\rho\kappa\Delta$ has no factor $1-1/b$.

\paragraph{Margins at calibration B.} The margins come in three groups.
\begin{enumerate}
\item \emph{Strict hypotheses} (Table~\ref{tab:margins}). These define the open parameter set. Each margin is the
  distance from the benchmark to the boundary of one hypothesis, with the exact cutoff in its formula.
\item \emph{Derived strict inequalities} (Table~\ref{tab:derived}). The proofs use these. They follow from the
  hypotheses, or they are exact or interval certificates at the benchmark auction.
\item \emph{Primitive domains and structural weak restrictions} (the list after the tables). These are not part of
  the open set of strict hypotheses.
\end{enumerate}
Every margin is positive. The smallest margins measure one distance: calibration B sets $M$ to the Laplace
manuscript's top posterior $e/(1+e)$, which lies $0.000334$ below $\tau(r_2)$.

\begin{table}[h]
\centering
\caption{Strict hypotheses and their margins at calibration B, prior $1/2$. Keys refer to
\texttt{registry/P1\_core.csv}.}\label{tab:margins}
\small
\begin{tabular}{@{}lllr@{}}
\toprule
Result & Hypothesis & Margin & Value at B\\
\midrule
Prop.~2 & $B_0(\pi)<c_H$ & $c_H-B_{r_0}(1/2)$ \hfill\texttt{margin\_weakcost\_B} & $1.195833$\\
 & $c_H<B_1(M)$ & $B_{r_1}(M)-c_H$ \hfill\texttt{margin\_recruit\_B} & $0.217155$\\
 & & $M-\tau(r_1)$ \hfill\texttt{margin\_recruit\_post\_B} & $0.026059$\\
 & $\Delta_0<k$ & $k-\Delta(r_0)$ \hfill\texttt{margin\_kweak\_B} & $0.003333$\\
 & $k<\rho\kappa\Delta_1$ & $\rho\kappa\Delta(r_1)-k$ \hfill\texttt{margin\_kstrong\_B} & $0.024824$\\
 & $B_2(M)<c_H$ & $c_H-B_{r_2}(M)$ \hfill\texttt{margin\_return\_B} & $0.002689$\\
 & & $\tau(r_2)-M$ \hfill\texttt{margin\_return\_post\_B} & $0.000334$\\
\midrule
Prop.~4 & $u>82/59$ & $u-82/59$ \hfill\texttt{margin\_u\_lo\_B} & $0.328451$\\
 & $u<2686/1559$ & $2686/1559-u$ \hfill\texttt{margin\_u\_hi\_B} & $0.004617$\\
 & $\lambda>\lambda_{\rm lo}(\beta)$ & $\lambda-\lambda_{\rm lo}(\beta)$ \hfill\texttt{margin\_lambda\_lo\_B} & $0.045583$\\
 & $\lambda<\lambda_{\rm hi}(\beta)$ & $\lambda_{\rm hi}(\beta)-\lambda$ \hfill\texttt{margin\_lambda\_hi\_B} & $0.000599$\\
 & $\pi>\pi_{\rm lo}$ & $1/2-\pi_{\rm lo}$ \hfill\texttt{margin\_prior\_lo\_B} & $0.032150$\\
 & $\pi<\pi_{\rm hi}$ & $\pi_{\rm hi}-1/2$ \hfill\texttt{margin\_prior\_hi\_B} & $0.000424$\\
\midrule
Prop.~5 & $\pi>\pi_A^{\rm lo}$ & $1/2-\pi_A^{\rm lo}$ \hfill\texttt{margin\_piA\_lo\_B} & $0.103825$\\
 & $\pi<\pi_A^{\rm hi}$ & $\pi_A^{\rm hi}-1/2$ \hfill\texttt{margin\_piA\_hi} & $0.140739$\\
\bottomrule
\end{tabular}
\end{table}

\begin{table}[h]
\centering
\caption{Derived strict inequalities used in the proofs, at the benchmark. The column ``Use'' names the step.
Exact rationals are \lab{analytical}; the other values are interval certificates (\lab{computer-assisted}).}
\label{tab:derived}
\small
\begin{tabular}{@{}llll@{}}
\toprule
Inequality & Margin and key & Value & Use\\
\midrule
$B_1(\pi)<c_H$ & $c_H-B_{r_1}(1/2)$ \hfill\texttt{margin\_implied\_prior\_r1} & $1.708333$ & implied by (H1), Prop.~2\\
$k<\rho\kappa\Delta_2$ & $\rho\kappa\Delta(r_2)-k$ \hfill\texttt{margin\_kstrong2\_B} & $0.043127$ & implied by (H2), Prop.~2\\
$B_2(\pi)<c_H$ & $c_H-B_{r_2}(1/2)$ \hfill\texttt{margin\_implied\_prior\_r2} & $1.865278$ & implied by (H3), Prop.~2\\
$u<19/3$ & $19/3-u$ \hfill\texttt{margin\_u\_kstrong\_B} & $4.615052$ & implied by the band, Prop.~4\\
\midrule
$\Delta>0$ on the path & $\Delta(r_0)$ \hfill\texttt{margin\_Delta\_pos\_B} & $0.016667$ & Lemma~\ref{lem:price}\\
$g_H>g_L$ on the path & $(g_H-g_L)(r_2)$ \hfill\texttt{margin\_gHgL\_B} & $8.061111$ & Lemma~\ref{lem:price}\\
$\tau(r_0)>1/2$ & $\tau(r_0)-1/2$ \hfill\texttt{margin\_tau\_r0\_half} & $0.133117$ & Prop.~3, Step 1\\
$J(r_0)<k$ & $k-J(r_0)$ \hfill\texttt{margin\_J\_r0\_B} & $0.013885$ & $r_0<r_A$\\
$\rho J(r_1)>k$ & $\rho J(r_1)-k$ \hfill\texttt{margin\_rhoJ\_r1\_B} & $0.029167$ & $r_J<r_1$\\
$\rho J(r_2)>k$ & $\rho J(r_2)-k$ \hfill\texttt{margin\_rhoJ\_r2\_B} & $0.043048$ & Lemma~\ref{lem:path}(b)\\
$\rho J(2)>k$ & $\rho J(2)-k$ \hfill\texttt{margin\_rhoJ\_2\_B} & $0.000982$ & $r_J<2$, Prop.~3, Step 7\\
$\tau(r_2)<3/4$ & $3/4-\tau(r_2)$ \hfill\texttt{margin\_tau\_r2\_3q\_B} & $0.018608$ & $M<3/4$, Lemma~\ref{lem:order}\\
$\tau(r_N)>1/2$ & $\tau(r_N)-1/2$ \hfill\texttt{margin\_tau\_rN\_half\_B} & $0.154463$ & $r_N<r_J$\\
$\tau_A<3/4$ & $3/4-\tau_A$ \hfill\texttt{margin\_tauA\_3q\_B} & $0.109261$ & Prop.~5, Step 3\\
$\tau_A<\pi_{\max}$ & $\pi_{\max}-\tau_A$ \hfill\texttt{margin\_pimax\_tauA\_B} & $0.009320$ & $r_A<r_*$, Prop.~5\\
$\Delta(r_*)[1-M(\pi_{\max})]>k$ & difference \hfill\texttt{margin\_cmend\_B} & $0.000343$ & Remark~\ref{rem:ties}(c) at B\\
\midrule
$r_0<r_A$ & $r_A-r_0$ \hfill\texttt{gap\_rA\_r0\_B} & $0.193958$ & Lemma~\ref{lem:order}\\
$r_A<r_F$ & $r_F-r_A$ \hfill\texttt{gap\_rF\_rA\_B} & $0.073169$ & Lemma~\ref{lem:order}\\
$r_F<r_N$ & $r_N-r_F$ \hfill\texttt{gap\_rN\_rF\_B} & $0.280750$ & Lemma~\ref{lem:order}\\
$r_N<r_J$ & $r_J-r_N$ \hfill\texttt{gap\_rJ\_rN\_B} & $0.218220$ & Lemma~\ref{lem:order}\\
$r_J<r_S$ & $r_S-r_J$ \hfill\texttt{gap\_rS\_rJ\_B} & $0.158051$ & Lemma~\ref{lem:order}\\
$r_S<r_1$ & $r_1-r_S$ \hfill\texttt{gap\_r1\_rS\_B} & $0.875852$ & Lemma~\ref{lem:order}\\
$r_1<r_C$ & $r_C-r_1$ \hfill\texttt{gap\_rC\_r1\_B} & $0.592659$ & Lemma~\ref{lem:order}\\
$r_C<r_2$ & $r_2-r_C$ \hfill\texttt{gap\_r2\_rC\_B} & $0.007341$ & Lemma~\ref{lem:order}\\
\bottomrule
\end{tabular}
\end{table}

Each gap encloses the whole difference of two roots: the script subtracts the full root brackets, so the
enclosure contains the true gap, not only a lower bound.

\paragraph{Primitive domains and structural weak restrictions.} These conditions define the model. They are not
part of the open set of strict hypotheses, and the margins table does not treat them as hypotheses.
\begin{itemize}
\item \emph{Primitive domains:} $0<\lambda<1$, $0<\beta<1/2$ (so that every liquidity order has positive
  probability), $0<\pi<1$, $0<\rho<1$, $k>0$, $c_H>0$, $c_L=0$, and $0<p<\ell<r<h$ on the path.
\item \emph{Structural weak restrictions:} $\beta\le1/3$ (the model's restriction; at calibration B its slack is
  $1/3-\beta=0.038678$, key \texttt{margin\_beta\_B}; only Lemma~\ref{lem:flow}(c) and Corollary~\ref{cor:ceiling}
  use it); $G_H$ dominates $G_L$; both supports lie above $p$; the facts (P1)--(P3).
\end{itemize}

\paragraph{Boundary equilibria at exact ties.} The strict-margin rule applies to hypotheses. It does not remove the
boundary points of an equilibrium set. Each point below is an exact equality, and the propositions keep it:
\begin{itemize}
\item $r=r_A$: the threshold pair T with $e^*=1$ (investor and challenger indifferent);
\item $r=r_F$: the full purchase F with $K_H=k$ (under IM);
\item $r=r_N$: no trade with $K_H=K_L=k$ (both classes are indifferent at prior $1/2$);
\item $r=r_J$: the threshold pair T with $e^*=\rho$ (under CM);
\item $r=r_S$: the full sale with $K_L=k$ (under IM);
\item $r=r_C$: the full purchase with a challenger tie (entry $\Ent_1$ under C1, any value in $[\rho,\Ent_1]$ under CM);
\item $\pi=\pi_A^{\rm lo}$: the singleton $C(\pi)=\{r_A\}$;
\item $\pi=\pi_{\max}$: the CM endpoint $r_*$ of Remark~\ref{rem:ties}.
\end{itemize}

%==========================================================================================
\section{Scripts, run logs and registry keys}\label{sec:scripts}

All scripts live in \texttt{audit/slot\_manuscript\_2026\_10\_23/scripts/P1\_core/}. They print their results and
write a run log beside themselves. Each one finishes in less than one minute on one process. Run them from the
repository root with \texttt{.venv/bin/python <script>}.
\begin{itemize}
\item \texttt{p1\_common.py} holds the shared objects. These are the benchmark primitives, the auction objects,
  exact polynomial arithmetic with Sturm sequences, outward intervals (\texttt{mpmath.iv}, 50 digits), the
  interval serialization and the registry writer. The serialization converts each interval endpoint exactly to a
  fraction. It prints the lower end rounded down and the upper end rounded up (24 places), so every printed
  enclosure contains its interval. No binary floating-point number enters a printed enclosure or a display.
\item \texttt{p1\_certificates.py} writes \texttt{p1\_certificates.log}. It checks the auction formulas by direct
  integration and gives the sign certificates of Lemma~\ref{lem:path}. It brackets $r_A$, $r_J$ and $r_*$ by
  exact bisection (width at most $10^{-40}$), with the signs at both ends, and it checks the order of the
  thresholds and the derived strict inequalities of Table~\ref{tab:derived}.
\item \texttt{p1\_numbers.py} writes \texttt{p1\_numbers.log} and \texttt{registry/P1\_core.csv} (147 rows). It
  computes every value at calibration B in this package, in exact fractions or in outward intervals. It checks
  that the interval for $e$ contains the series bracket $S_{110}<e<S_{110}+112/(111\cdot111!)$. A display counts as
  certified only if both exact ends of its enclosure round to it. The script flags every other case, and it flags
  none. Each gap between two thresholds subtracts the whole root brackets.
\item \texttt{p1\_enumerate.py} writes \texttt{p1\_enumerate.log} (\lab{numerical diagnostic}). It rebuilds the
  equilibrium set from the flow law at 241 strengths, for four pairs of tie rules and two calibrations (1,928
  cases), and finds no difference from Proposition~\ref{prop:3}. It also scans coexistence at 13 priors, 481
  strengths and four rule pairs in both calibrations (50,024 cases), and finds no difference from
  Proposition~\ref{prop:5} and Remark~\ref{rem:ties}. At a zero rate the posterior equals the prior, so these
  comparisons use exact fractions; every comparison is decided.
\item \texttt{p1\_compare.py} writes \texttt{p1\_compare.log}. It compares the registry with its first version
  (no display changes, no removed keys) and with the independent exact intervals of the check: each of the 70
  shared enclosures contains the independent interval.
\end{itemize}
The registry uses the starter keys of the brief. It adds keys in the same pattern, for example \texttt{rJ},
\texttt{rS\_B}, \texttt{rstar}, \texttt{piA\_lo\_B}, \texttt{piA\_hi}, \texttt{pimax}, \texttt{entry\_rA\_B},
\texttt{entry\_mb\_r140\_B} and \texttt{entry\_th\_r155\_B}. Tables~\ref{tab:margins} and~\ref{tab:derived} give the
keys of the margins.

%==========================================================================================
\section{Differences from the brief and the earlier memos}\label{sec:differences}

\begin{enumerate}
\item \emph{The value of $r_F$ at calibration B.} The earlier study reports $r_F=1.460985$. That number is
  $\fr(4/55)$, the value at calibration A, where $M=29/40$; it is correct there. At calibration B,
  $r_F=\fr(k(1+e))=1.467127$, the value that the brief now states.
\item \emph{Coexistence at other priors.} The model page of the earlier study states that outside the window
  $(0.396,0.641)$ coexistence can still occur but starts later. Proposition~\ref{prop:5} gives the exact set under
  IM and C1: it starts later only for $\pi\in[0.640739,0.650059)$. For $\pi<0.396175$, and for $\pi\ge0.650059$, no
  strength on the path has coexistence under C1. Under CM the endpoint $\pi=\pi_{\max}$ adds the single strength
  $r_*$ (Remark~\ref{rem:ties}).
\item \emph{The investor rule I0.} The earlier memos give the coexistence interval $(r_F,r_N]$ under I0. Under I0
  the equilibrium set is also empty on $(r_N,r_S]$, that is $(1.747878,2.124148]$ at B. The low class has no
  admissible rate there (Lemma~\ref{lem:L}).
\item \emph{Where $\beta\le1/3$ enters.} It enters only the bound at the zero order (Lemma~\ref{lem:flow}(c)) and
  Corollary~\ref{cor:ceiling}. Propositions 2--5 hold for every $\beta<1/2$.
\item \emph{Entry at $r_1$.} With the declared $\beta=0.294655$, entry at $r_1$ is $0.5227569$; the unrounded
  $\beta$ of the earlier study gives $0.5227573$. Both display as $0.522757$.
\item \emph{The threshold pair at B.} The earlier values $0.350846$, $0.328628$ and $0.310872$ of entry on T at
  $1.55$, $1.60$ and $1.65$ belong to calibration A. At calibration B they are $0.341462$, $0.321311$ and $0.305208$.
\item \emph{The meaning of $u$.} The first version of this package called $u$ the odds that an order in one
  direction is informed. That reading holds for no strategy pair in general: under full orders the odds are
  $\pi u$ after a purchase and $(1-\pi)u$ after a sale. The number $u$ is the informed presence mass relative to the
  liquidity mass in one direction, and $1+u$ is the largest likelihood ratio of an order (Section~\ref{sec:model}).
\end{enumerate}
The package confirms these earlier results:
\begin{itemize}
\item the corrected hypotheses (H1)--(H3) and the ownership condition;
\item $r_A=1.393958$, $r_N=1.747878$ and $r_C=3.592659$;
\item the prior range $(\pi_{\rm lo},\pi_{\rm hi})\approx(0.467850,0.500424)$;
\item entry $0.522757$ and ownership $0.324401$ at $r_1$;
\item the return margin $0.000334$ in posterior units.
\end{itemize}

\end{document}
